\documentclass[12pt]{book} % Change to book
\usepackage{amsmath} % for mathematical expressions
\usepackage{xcolor} % for colored text
\usepackage{graphicx} % for including images
\usepackage{url}
\usepackage{booktabs} % for better-looking tables
\usepackage{makecell}
\usepackage{indentfirst}
\usepackage{algorithm}
\usepackage{algpseudocode}
\usepackage{amssymb}
\usepackage{placeins} % For \FloatBarrier
\usepackage{listings}
% Define the language style for Python code
\lstset{
    language=Python,
    basicstyle=\ttfamily,
    keywordstyle=\color{blue},
    stringstyle=\color{red},
    commentstyle=\color{gray},
    morecomment=[l][\color{magenta}]{\#},
    frame=single,
    breaklines=true
}
\begin{document}
\chapter{Graph Algorithms}
\section{Introduction to Graph Theory}

Graph theory is a vibrant area of mathematics that studies graphs—simple structures that model relationships between pairs of objects. These objects are represented as vertices (or nodes), and the connections between them are edges (or links). Whether in computer science, biology, or social science, graph theory plays a crucial role by providing a framework to solve problems involving connected data.


Graphs come in several flavors:
- \textbf{Undirected graphs} where edges have no direction.
- \textbf{Directed graphs} where edges point from one vertex to another.
- \textbf{Weighted graphs} where edges carry weights, representing costs, lengths, or capacities.

\textbf{Exploring Graph Theory}

Studying graph theory involves:
- Examining the properties of graphs to understand their structure.
- Developing algorithms to tackle problems like finding the shortest path between nodes or determining the most efficient way to connect various points.
- Applying graph theory concepts to practical scenarios across different fields, showing how these abstract ideas help solve real-world problems.

In summary, graph theory not only enriches our mathematical understanding but also enhances our ability to address complex issues in various scientific and practical domains.


\paragraph{Algorithmic Example: Breadth-First Search (BFS)}
\FloatBarrier
Breadth-First Search is a fundamental graph traversal algorithm used to explore the vertices of a graph level by level. It starts at a given vertex and visits its neighbors before moving on to the next level of neighbors.
\FloatBarrier
The BFS algorithm can be described as follows:
\FloatBarrier
\begin{algorithm}
\caption{Breadth-First Search}
\begin{algorithmic}
\State Initialize a queue $Q$ with the starting vertex
\State Mark the starting vertex as visited
\While {$Q$ is not empty}
    \State Dequeue a vertex $v$ from $Q$
    \For {each neighbor $u$ of $v$}
        \If {$u$ is not visited}
            \State Mark $u$ as visited and enqueue $u$ into $Q$
        \EndIf
    \EndFor
\EndWhile
\end{algorithmic}
\end{algorithm}
\FloatBarrier
The Python code for Breadth-First Search algorithm can be implemented as follows:
\FloatBarrier
\begin{lstlisting}
def bfs(graph, start):
    visited = set()
    queue = [start]
    visited.add(start)
    
    while queue:
        node = queue.pop(0)
        for neighbor in graph[node]:
            if neighbor not in visited:
                visited.add(neighbor)
                queue.append(neighbor)
\end{lstlisting}


\subsection{Definition of Graphs}

At its core, a graph \( G = (V, E) \) in graph theory is made up of vertices \( V \) and edges \( E \), where each edge connects a pair of vertices. This setup is used to model the relationships between different entities. Depending on the nature of these relationships, edges can be directed or undirected, and the graph can either carry weights on its edges or not.
\newline
\newline
{\textbf{Exploring Different Types of Graphs}}


\FloatBarrier
- \textbf{Undirected Graphs}: Here, edges are bidirectional, meaning there's no specific direction to the connection. They simply link two vertices together, allowing travel in both directions.


\paragraph{Algorithmic Example - Breadth-First Search (BFS):}
Breadth-First Search is a graph traversal algorithm that visits all the vertices in a graph in breadth-first order starting from a chosen vertex. Here is the algorithm:

\begin{algorithm}
\caption{Breadth-First Search (BFS)}
\begin{algorithmic}
\State Let $G = (V, E)$ be the graph, $s$ be the start vertex
\State Initialize an empty queue $Q$
\State Enqueue the start vertex $s$ into $Q$
\While{$Q$ is not empty}
    \State Dequeue a vertex $v$ from $Q$
    \For{each neighbor $u$ of $v$}
        \If{$u$ is not visited}
            \State Mark $u$ as visited
            \State Enqueue $u$ into $Q$
        \EndIf
    \EndFor
\EndWhile
\end{algorithmic}
\end{algorithm}

The Python code for Breadth-First Search (BFS) algorithm can be implemented as follows:

\begin{lstlisting}
from collections import deque

def bfs(graph, start):
    visited = set()
    queue = deque([start])
    visited.add(start)

    while queue:
        node = queue.popleft()
        print(node)

        for neighbor in graph[node]:
            if neighbor not in visited:
                visited.add(neighbor)
                queue.append(neighbor)

# Example Usage
graph = {0: [1, 2], 1: [2], 2: [0, 3], 3: [3]}
bfs(graph, 2)
\end{lstlisting}

\subsection{Types of Graphs}

Graph theory classifies graphs into several types based on their properties and the nature of the connections between their nodes. Understanding these types helps in applying the right kind of graph to solve specific problems effectively.

1. \textbf{Directed Graphs}: In these graphs, edges point from one vertex to another, establishing a clear direction from the source to the destination. This structure is useful in scenarios like city traffic flows where roads have definite directions.

2. \textbf{Undirected Graphs}: These graphs feature edges that don’t have a directed flow, meaning the relationship between the connected vertices is bidirectional. They are ideal for modeling undirected connections like telephone lines between cities.

3. \textbf{Weighted Graphs}: Here, edges carry weights, which could be costs, distances, or any other measurable attribute. This type is particularly useful in route planning and network optimization, where the cost of traversing edges varies significantly.

4. \textbf{Connected Graphs}: Every pair of vertices in a connected graph is linked by a path. This type ensures there are no isolated vertices or disjoint subgraphs, which is crucial for network design to ensure every node is reachable.

5. \textbf{Complete Graphs}: These graphs feature an edge between every pair of vertices. They represent a fully connected network, often used in problems requiring exhaustive interconnection, like circuit design or social network analysis.

Each type of graph has its specific use cases, and choosing the correct type is key to efficiently solving complex graph-based problems.


\subsubsection{Example - Depth-First Search Algorithm}

Let's consider a simple example of implementing the Depth-First Search (DFS) algorithm on a graph. DFS is used to traverse or search a graph in a specific manner.
\FloatBarrier
Algorithm for Depth-First Search:
\FloatBarrier
\begin{algorithm}
\caption{Depth-First Search}
\begin{algorithmic}[1]
\Function{DFS}{graph, node}
\State Mark node as visited
\For {each neighbor of node}
\If {neighbor is not visited}
\State \Call{DFS}{graph, neighbor}
\EndIf
\EndFor
\EndFunction
\end{algorithmic}
\end{algorithm}
\FloatBarrier
Now, let's implement the equivalent Python code for the DFS algorithm:
\FloatBarrier
\begin{lstlisting}
\begin{lstlisting}
\begin{algorithmic}
\Function{DFS}{graph, node}
\State visited[node] = True
\For neighbor in graph[node]:
\If not visited[neighbor]:
\State \Call{DFS}{graph, neighbor}
\EndIf
\EndFor
\EndFunction
\end{algorithmic}
\end{lstlisting}
\FloatBarrier
\subsubsection{Directed vs. Undirected Graphs}
In graph theory, graphs can be categorized into two main types based on the nature of their edges - directed graphs and undirected graphs.
\FloatBarrier
\textbf{Directed Graphs}:
\FloatBarrier
In a directed graph, the edges have a direction associated with them. This means that the relationship between any two nodes is asymmetric. If there is an edge from node A to node B, it does not imply the existence of an edge from node B to node A.
\FloatBarrier
\paragraph{Example: Depth-First Search (DFS) on a Directed Graph}
\FloatBarrier
\begin{algorithm}
\caption{Depth-First Search (DFS) on a Directed Graph}
\begin{algorithmic}[1]
\Function{DFS}{$G, v$}
    \State Mark vertex $v$ as visited
    \For{each adjacent vertex $u$ of $v$}
        \If{$u$ is not visited}
            \State \textsc{DFS}($G, u$)
        \EndIf
    \EndFor
\EndFunction
\end{algorithmic}
\end{algorithm}
\FloatBarrier
\begin{verbatim}
def DFS(G, v, visited):
    visited[v] = True
    print("Visited:", v)
    for u in G[v]:
        if not visited[u]:
            DFS(G, u, visited)

# Example usage:
# Define the directed graph as an adjacency list
G = {
    0: [1, 2],
    1: [2],
    2: [0, 3],
    3: [3]
}

# Initialize visited array
visited = [False] * len(G)

# Perform DFS starting from vertex 2
DFS(G, 2, visited)
\end{verbatim}
\FloatBarrier

\textbf{Algorithm Explanation:}

The Depth-First Search (DFS) algorithm is used to traverse or search a graph. In the case of a directed graph, DFS starts from a given vertex $v$ and explores as far as possible along each branch before backtracking.

\textbf{Mathematical Detail:}

Let $G(V, E)$ be a directed graph with vertices $V$ and edges $E$. The DFS algorithm visits each vertex in $G$ exactly once, making it $O(|V|)$ in time complexity, where $|V|$ is the number of vertices in $G$. 

Suppose the DFS algorithm starts from vertex $v_0$. It explores all vertices reachable from $v_0$ in a depth-first manner. Let $n$ be the number of vertices reachable from $v_0$. The DFS algorithm has a space complexity of $O(n)$ because it stores information about the vertices it has visited, typically using a stack or recursion.

The DFS algorithm can be used for various purposes in directed graphs, such as detecting cycles, finding strongly connected components, and topological sorting.
\FloatBarrier
\textbf{Undirected Graphs}:
\FloatBarrier
In an undirected graph, the edges do not have any direction associated with them. This means that the relationship between any two nodes is symmetrical. If there is an edge between node A and node B, it implies the existence of an edge between node B and node A.
\FloatBarrier
\paragraph{Example: Depth-First Search (DFS) on a Undirected Graph}
\FloatBarrier
\begin{verbatim}
def DFS(G, v, visited):
    visited[v] = True
    print("Visited:", v)
    for u in G[v]:
        if not visited[u]:
            DFS(G, u, visited)

# Example usage:
# Define the directed graph as an adjacency list
G = {
    0: [1, 2],
    1: [2],
    2: [0, 3],
    3: [3]
}

# Initialize visited array
visited = [False] * len(G)

# Perform DFS starting from vertex 2
DFS(G, 2, visited)
\end{verbatim}
\FloatBarrier
\paragraph{Python Code:}
\FloatBarrier
\begin{verbatim}
def DFS(G, v, visited):
    visited[v] = True
    print("Visited:", v)
    for u in G[v]:
        if not visited[u]:
            DFS(G, u, visited)

# Example usage:
# Define the undirected graph as an adjacency list
G = {
    0: [1, 2],
    1: [0, 2],
    2: [0, 1, 3],
    3: [2]
}

# Initialize visited array
visited = [False] * len(G)

# Perform DFS starting from vertex 0
DFS(G, 0, visited)
\end{verbatim}
\FloatBarrier
\subsubsection{Weighted vs. Unweighted Graphs}

In graph theory, graphs are generally divided into two categories based on edge characteristics: weighted and unweighted.

\textbf{Unweighted Graphs}: These graphs treat all edges equally, as each edge has the same weight or no weight at all. Unweighted graphs are ideal for situations where only the connection matters, not the strength or capacity of that connection. They simplify modeling relationships where every link is of equal value, such as friendships in a social network where the presence of a connection is more important than its intensity.

\textbf{Weighted Graphs}: In contrast, weighted graphs assign a numerical weight to each edge, which can represent distance, cost, time, or other metrics. This feature makes weighted graphs essential for accurately modeling real-world problems where not all connections are created equal. For instance, in a road network, edges can represent the distance or travel time between locations, affecting route planning.

\textbf{Implications for Algorithms}: The type of graph impacts how algorithms perform. For example, finding the shortest path in a network is straightforward in unweighted graphs where each step counts the same. However, in weighted graphs, algorithms like Dijkstra's need to consider the weights to optimize the path based on actual travel costs or distances.

Understanding whether to use a weighted or unweighted graph depends on the specific requirements of the application and the nature of the entities being modeled. This decision is crucial as it influences both the complexity of the problem and the choice of algorithms for processing the graph.


\subsubsection*{Algorithmic Example: Dijkstra's Algorithm}

Dijkstra's algorithm is a widely-used algorithm for finding the shortest path between nodes in a weighted graph. Given a weighted graph, the algorithm calculates the shortest path from a source node to all other nodes in the graph.
\FloatBarrier
Let's consider a weighted graph represented by an adjacency matrix:
\FloatBarrier
\[ \begin{pmatrix}
0 & 4 & 0 & 0 & 0 \\
0 & 0 & 8 & 0 & 0 \\
0 & 0 & 0 & 7 & 0 \\
0 & 0 & 0 & 0 & 5 \\
0 & 0 & 0 & 0 & 0
\end{pmatrix} \]
\FloatBarrier
The Python implementation of Dijkstra's algorithm for finding the shortest path from a given source node to all other nodes in the graph can be represented as follows:
\FloatBarrier
\begin{algorithm}
\caption{Dijkstra's Algorithm}
\begin{algorithmic}[1]
\Procedure{Dijkstra}{graph, source}
\State $distances \gets$ Array filled with $\infty$ of size $|V|$ \Comment{Initialize distances array}
\State $distances[source] \gets 0$ \Comment{Distance from source to itself is 0}
\State $visited \gets$ Empty set \Comment{Initialize set of visited nodes}

\While{There are unvisited nodes}
    \State $v \gets$ Node with the minimum distance in $distances$ \Comment{Select unvisited node}
    \State $visited$.add($v$) \Comment{Mark node as visited}
    
    \For{each neighbor $u$ of $v$}
        \State $alt \gets$ distance[$v$] + weight of edge $(v, u)$
        \If{$alt < distances[u]$}
            \State $distances[u] \gets alt$ \Comment{Update distance if shorter path found}
        \EndIf
    \EndFor
\EndWhile

\State \textbf{return} $distances$
\EndProcedure
\end{algorithmic}
\end{algorithm}
\FloatBarrier
This algorithm calculates the shortest path from the given source node to all other nodes in the graph. The resulting distances array will contain the shortest path distances from the source node to each node in the graph.

\subsection{Representations of Graphs}

In computer science, how we represent graphs in data structures is crucial for efficient storage and manipulation. Two of the most common graph representations are the adjacency matrix and the adjacency list.

\textbf{Adjacency Matrix}: This representation uses a 2D array where both rows and columns correspond to vertices in the graph. An entry at \(A[i][j]\) indicates the presence of an edge between vertex \(i\) and vertex \(j\). If an edge exists, the entry is set to 1 or to the weight of the edge in a weighted graph. If no edge exists, it's set to 0, or \(\infty\) in weighted graphs to denote an absence of direct paths.

\textbf{Adjacency List}: Alternatively, an adjacency list uses a list of lists. Each list corresponds to a vertex and contains the vertices that are directly connected to it. This method is space-efficient, especially for sparse graphs, as it only stores information about existing edges.

Each representation has its advantages depending on the operations required and the density of the graph. The adjacency matrix makes it quick and easy to check if an edge exists between any two vertices but can be space-intensive for large sparse graphs. On the other hand, adjacency lists are more space-efficient for sparse graphs but can require more time to check for the existence of a specific edge. Choosing the right representation is key to optimizing the performance of graph algorithms.


Algorithm Example: Finding the degree of a vertex in a graph using an adjacency list representation.
\FloatBarrier
\begin{algorithm}
\caption{Find Degree of a Vertex using Adjacency List}
\begin{algorithmic}
\State \textbf{Input:} Graph $G$ represented as an adjacency list, Vertex $v$
\State \textbf{Output:} Degree of Vertex $v$

\Function{Degree}{$G, v$}
    \State $degree \gets |G[v]|$ \Comment{Get the length of the adjacency list for vertex $v$}
    \State \textbf{return} $degree$
\EndFunction
\end{algorithmic}
\end{algorithm}
\FloatBarrier
Equivalent Python code:
\FloatBarrier
\begin{lstlisting}
\begin{algorithmic}
\Function{Degree}{G, v}
    \State degree = len(G[v])  \Comment{Get the length of the adjacency list for vertex v}
    \State return degree
\EndFunction
\end{algorithmic}
\end{lstlisting}
\FloatBarrier
These representations have their own advantages and are selected based on the operations expected to be performed on the graph data structure. The choice of representation can significantly impact the efficiency of algorithms that work on graphs.

\subsubsection{Adjacency Matrix}

An adjacency matrix is a way to represent a graph as a square matrix where the rows and columns are indexed by vertices. The entry at row $i$ and column $j$ represents whether there is an edge from vertex $i$ to vertex $j$. If the graph is unweighted, the entry will be either 0 (no edge) or 1 (edge exists). In the case of a weighted graph, the entry can hold the weight of the edge.
\FloatBarrier
Algorithmic Example:
\FloatBarrier
Let's consider a simple undirected graph with 4 vertices and the following edges: (1, 2), (2, 3), (3, 4), (4, 1).
\FloatBarrier
The adjacency matrix for this graph would look like:
\FloatBarrier
\[
\begin{bmatrix}
0 & 1 & 0 & 1 \\
1 & 0 & 1 & 0 \\
0 & 1 & 0 & 1 \\
1 & 0 & 1 & 0 \\
\end{bmatrix}
\]
\FloatBarrier
This matrix reflects the connections between vertices in the graph.
\FloatBarrier
Equivalent Python code for constructing the adjacency matrix:
\FloatBarrier
\begin{lstlisting}
def adjacency_matrix(graph):
    n = len(graph)
    adj_matrix = [[0] * n for _ in range(n)]
    
    for edge in graph:
        node1, node2 = edge
        adj_matrix[node1][node2] = 1
        adj_matrix[node2][node1] = 1
    
    return adj_matrix

graph = [(0, 1), (1, 2), (2, 3), (3, 0)]
adj_matrix = adjacency_matrix(graph)
print(adj_matrix)
\end{lstlisting}
\FloatBarrier
This Python function takes a graph as an input (list of edges) and constructs the corresponding adjacency matrix. In the example provided, the graph has the edges (0, 1), (1, 2), (2, 3), and (3, 0). The resulting adjacency matrix is printed out.
Sure, I can help you with that. Here is an expanded explanation of the \texttt{Adjacency List} approach in LaTeX format along with an algorithmic example and its equivalent Python code.

\subsubsection{Adjacency List}

In the adjacency list representation of a graph, each vertex in the graph is associated with a list of its neighboring vertices. This representation is quite efficient for sparse graphs where the number of edges is much smaller compared to the number of possible edges. It requires less memory storage compared to the adjacency matrix representation, especially when the graph is sparse.

When using an adjacency list to represent a graph, we can easily find all the neighbors of a specific vertex by looking at its associated list. This makes it efficient for algorithms like traversals and shortest-path algorithms.
\FloatBarrier
Algorithm:
\FloatBarrier
Let \( G(V, E) \) be a graph with vertices \( V = \{v_1, v_2, ..., v_n\} \) and edges \( E = \{e_1, e_2, ..., e_m\} \). The adjacency list representation of \( G \) is a set of \( n \) lists, one for each vertex \( v_i \), where each list contains the vertices adjacent to \( v_i \).
\FloatBarrier
\textbf{Example:}
\FloatBarrier
Let's consider a simple undirected graph with 4 vertices and 4 edges:
\FloatBarrier
\[ V = \{v_1, v_2, v_3, v_4\} \]
\[ E = \{(v_1, v_2), (v_2, v_3), (v_3, v_4), (v_1, v_4)\} \]
\FloatBarrier
The adjacency list for this graph would be:
\FloatBarrier
\[ v_1: [v_2, v_4] \]
\[ v_2: [v_1, v_3] \]
\[ v_3: [v_2, v_4] \]
\[ v_4: [v_1, v_3] \]
\FloatBarrier
\textbf{Algorithmic example:}
\begin{algorithm}
\caption{Create Adjacency List of a Graph}
\begin{algorithmic}
\State Let $adjList$ be an empty dictionary
\For{each edge $(u, v)$ in $E$}
    \State Add $v$ to $adjList[u]$
    \State Add $u$ to $adjList[v]$
\EndFor
\end{algorithmic}
\end{algorithm}
\FloatBarrier
\textbf{Python Code:}
\FloatBarrier
\begin{verbatim}
# Initialize an empty dictionary for the adjacency list
adjList = {}

# Loop through each edge (u, v) in the set of edges E
for u, v in E:
    # Add v to the adjacency list of u
    if u in adjList:
        adjList[u].append(v)
    else:
        adjList[u] = [v]
    
    # Add u to the adjacency list of v
    if v in adjList:
        adjList[v].append(u)
    else:
        adjList[v] = [u]
\end{verbatim}
\FloatBarrier

\section{Graph Traversal Techniques}
Graph traversal is the process of visiting all the nodes in a graph in a systematic way. There are several techniques for graph traversal, some of the common ones include:
\begin{itemize}
    \item Breadth-first Search (BFS): In BFS, we visit all the neighbors of a node before moving on to the next level of nodes.
    \item Depth-first Search (DFS): In DFS, we explore as far as possible along each branch before backtracking.
\end{itemize}
\FloatBarrier
\textbf{Algorithmic Example: Breadth-first Search (BFS)}
\FloatBarrier
   \begin{algorithm}
     \caption{BFS Algorithm}
     \begin{algorithmic}[1]
       \State Initialize a queue $Q$ with the starting node $s$
       \State Mark $s$ as visited
       \While{$Q$ is not empty}
         \State Dequeue a node $v$ from $Q$
         \For{each neighbor $w$ of $v$}
           \If{$w$ is not visited}
             \State Mark $w$ as visited
             \State Enqueue $w$ in $Q$
           \EndIf
         \EndFor
       \EndWhile
     \end{algorithmic}
   \end{algorithm}
\FloatBarrier
\textbf{Python Code Equivalent for BFS Algorithm}
\FloatBarrier
   \begin{lstlisting}
     def bfs(graph, start):
         visited = set()
         queue = [start]
         visited.add(start)
         while queue:
             node = queue.pop(0)
             for neighbor in graph[node]:
                 if neighbor not in visited:
                     visited.add(neighbor)
                     queue.append(neighbor)
   \end{lstlisting}
   \FloatBarrier

\subsection{Breadth-First Search (BFS)}
Breadth-First Search (BFS) is an algorithm used for traversing or searching tree or graph data structures. It starts at the tree root or any arbitrary node of a graph, and explores all of the neighbor nodes at the present depth prior to moving on to the nodes at the next depth level. This ensures that the algorithm traverses the graph in layers. BFS is often used to find the shortest path in unweighted graphs.

\subsubsection{Algorithm Overview}


\FloatBarrier
\begin{algorithm}
\caption{Breadth-First Search (BFS)}
\begin{algorithmic}
\State $Q \gets$ Queue data structure
\State $visited \gets$ Set to keep track of visited nodes
\State $Q.push(start)$ \Comment{Start with the initial node}
\State $visited.add(start)$
\While{$Q$ is not empty}
    \State $current \gets Q.pop()$
    \For{each neighbor $n$ of $current$}
        \If{$n$ is not in $visited$}
            \State $Q.push(n)$
            \State $visited.add(n)$
        \EndIf
    \EndFor
\EndWhile
\end{algorithmic}
\end{algorithm}
\FloatBarrier
\textbf{Python Code Equivalent:}
\FloatBarrier

\begin{lstlisting}
from collections import deque

def bfs(graph, start):
    visited = set()
    q = deque()
    
    q.append(start)
    visited.add(start)
    
    while q:
        current = q.popleft()
        for neighbor in graph[current]:
            if neighbor not in visited:
                q.append(neighbor)
                visited.add(neighbor)
    
    return visited
\end{lstlisting}
\FloatBarrier

\subsubsection{Applications and Examples}

The Breadth First Search (BFS) algorithm is a cornerstone of graph traversal techniques, widely appreciated for its simplicity and effectiveness across various fields. Here are some key applications:

\textbf{Shortest Path and Distance Calculation}
Primarily, BFS is used to identify the shortest path and calculate distances in unweighted graphs. It explores the graph layer by layer from the source node, ensuring that the shortest paths are identified efficiently.

\textbf{Example:} In navigation systems, BFS helps calculate the most direct route between locations, considering each junction or road one by one from the starting point.

\textbf{Connected Components}
BFS is also instrumental in detecting connected components in undirected graphs. By launching BFS from various nodes, it's possible to group nodes into clusters based on connectivity.

\textbf{Example:} In social network analysis, BFS can determine groups of users who are interconnected, providing insights into community structures.

\textbf{Bipartite Graph Detection}
Another application of BFS is in verifying whether a graph is bipartite, meaning it can be split into two groups where edges only connect nodes from opposite groups.

\textbf{Example:} This feature of BFS is useful in scheduling tasks, like in schools or universities, to ensure that courses or exams can be arranged without overlaps in resources or time slots.

\textbf{Network Broadcast and Message Propagation}
In network protocols, BFS facilitates the efficient propagation of messages or data packets across all reachable parts of the network, ensuring minimal delay and redundancy.

\textbf{Example:} BFS-based algorithms in Ethernet networks help in broadcasting messages to all devices, ensuring each one receives the data in the shortest possible time.

These examples highlight the versatility of BFS in solving practical problems efficiently, from route planning and community detection to network communications and scheduling.


\subsection{Depth-First Search (DFS)}

Depth First Search (DFS) is a graph traversal algorithm that explores as far as possible along each branch before backtracking. It starts at a selected node (often called the "root" node) and explores as far as possible along each branch before backtracking. DFS uses a stack to keep track of the nodes to visit next.

\subsubsection{Algorithm Overview}

The DFS algorithm can be described as follows:
\FloatBarrier
\begin{itemize}
    \item Choose a starting node and mark it as visited.
    \item Explore each adjacent node that has not been visited.
    \item Repeat the process recursively for each unvisited adjacent node.
    \item Backtrack when there are no more unvisited adjacent nodes.
\end{itemize}
\FloatBarrier
\textbf{Pseudocode}
\FloatBarrier
The DFS algorithm can be implemented using the following pseudocode:
\FloatBarrier
\begin{algorithm}
\caption{Depth First Search (DFS)}
\begin{algorithmic}[1]
\Function{DFS}{$G, start$}
    \State $visited \gets \emptyset$
    \State \Call{DFS\_Visit}{$G, start, visited$}
\EndFunction

\Function{DFS\_Visit}{$G, node, visited$}
    \State $visited.\text{add}(node)$
    \For{$\text{neighbor}$ \textbf{in} $G[node]$}
        \If{$\text{neighbor} \notin visited$}
            \State \Call{DFS\_Visit}{$G, neighbor, visited$}
        \EndIf
    \EndFor
\EndFunction
\end{algorithmic}
\end{algorithm}
\FloatBarrier
In the above pseudocode, $G$ represents the graph, $start$ is the starting node, and $visited$ is a set to keep track of visited nodes.
\FloatBarrier

\textbf{Complexity Analysis}
\FloatBarrier
The time complexity of the DFS algorithm is $O(V + E)$, where $V$ is the number of vertices and $E$ is the number of edges in the graph.

\FloatBarrier
\textbf{Python Code Equivalent:}
\FloatBarrier
\begin{algorithm}
\caption{Depth First Search (DFS)}
\begin{algorithmic}[1]
\Function{DFS}{$G, start$}
    \State $visited \gets \emptyset$
    \State \Call{DFS\_Visit}{$G, start, visited$}
\EndFunction

\Function{DFS\_Visit}{$G, node, visited$}
    \State $visited.\text{add}(node)$
    \For{$\text{neighbor}$ \textbf{in} $G[node]$}
        \If{$\text{neighbor} \notin visited$}
            \State \Call{DFS\_Visit}{$G, neighbor, visited$}
        \EndIf
    \EndFor
\EndFunction
\end{algorithmic}
\end{algorithm}
\FloatBarrier

\subsubsection{Applications and Examples}

Depth First Search (DFS) is a versatile graph traversal technique used extensively across various fields for its depth-focused exploration strategy. Here are some notable applications:

\textbf{Graph Traversal}
DFS excels in navigating through complex structures, diving deep into each path before backtracking. This method is ideal for searching through large networks like social media sites, the internet, or interconnected computer systems.

\textbf{Cycle Detection}
In cycle detection, DFS proves useful by identifying cycles within graphs. If DFS revisits a node via a back edge (an edge connecting to an ancestor in the DFS tree), it signals a cycle, crucial in many applications such as circuit testing or workflow analysis.

\textbf{Topological Sorting}
DFS facilitates topological sorting in directed acyclic graphs (DAGs). It orders nodes linearly ensuring that for any directed edge \(uv\), node \(u\) precedes \(v\). This is particularly useful in scheduling tasks, compiling data, or course prerequisite planning.

\textbf{Connected Components}
In undirected graphs, DFS helps identify connected components, enabling analysis of network clusters or related groups in data.

\textbf{Maze Solving}
DFS's backtracking feature makes it ideal for maze-solving applications, exploring all potential paths until an exit is found, often used in gaming and robotics simulations.

\textbf{Strongly Connected Components}
For directed graphs, DFS identifies strongly connected components, ensuring that every vertex is reachable from any other vertex, which is crucial for understanding the robustness of networks.

\textbf{Path Finding}
DFS is effective in finding paths between vertices in a graph, exploring all possible routes from a start to an endpoint, useful in routing and navigation systems.

\textbf{Network Analysis}
In network analysis, DFS is instrumental in identifying critical structures like bridges and articulation points, which help in assessing network vulnerability and stability.

\subsection{Comparing BFS and DFS}

While both Breadth-First Search (BFS) and Depth-First Search (DFS) are foundational for graph traversal, they differ significantly in their approach and applications. BFS explores the graph level by level using a queue, making it suitable for shortest path calculations and ensuring all nodes at the current depth are explored before moving deeper. DFS, on the other hand, dives deep into the graph using a stack or recursion, which is useful for tasks that need to explore all paths like puzzle solving or cycle detection. Choosing between BFS and DFS depends on the specific requirements of the problem at hand.
\newline
Let's compare the two algorithms using an example:
\FloatBarrier
\begin{algorithm}
\caption{Comparison of BFS and DFS}
\begin{algorithmic}[1]
\State Let $G(V, E)$ be the graph with vertex set $V$ and edge set $E$
\State Let $s$ be the starting vertex for traversal
\State Initialize an empty queue $Q$
\State Initialize an empty stack $S$
\State Enqueue vertex $s$ into $Q$ and push it onto $S$

\While{$Q$ is not empty}
    \State Dequeue a vertex $v$ from $Q$ and pop $v$ from $S$
    \For{each neighbor $u$ of $v$}
        \If{vertex $u$ is unvisited}
            \State Enqueue $u$ into $Q$ and push it onto $S$
        \EndIf
    \EndFor
\EndWhile
\end{algorithmic}
\end{algorithm}
\FloatBarrier
In this algorithmic example, we demonstrate how BFS and DFS work based on the exploration of neighboring vertices starting from a given vertex. The Python code equivalent for this comparison is as follows:
\FloatBarrier
\begin{lstlisting}
# Python code for BFS and DFS comparison
from collections import deque

def bfs(graph, start):
    visited = set()
    queue = deque([start])
    while queue:
        vertex = queue.popleft()
        if vertex not in visited:
            visited.add(vertex)
            queue.extend(graph[vertex] - visited)
    return visited

def dfs(graph, start):
    visited = set()
    stack = [start]
    while stack:
        vertex = stack.pop()
        if vertex not in visited:
            visited.add(vertex)
            stack.extend(graph[vertex] - visited)
    return visited

# Example usage
graph = {1: {2, 3}, 2: {4, 5}, 3: {6}, 4: set(), 5: {7}, 6: set(), 7: set()}
print("BFS traversal:", bfs(graph, 1))
print("DFS traversal:", dfs(graph, 1))
\end{lstlisting}
\FloatBarrier
These implementations show how BFS and DFS can be used to traverse a graph and explore its vertices. Both algorithms offer unique approaches to graph traversal with different applications and efficiencies.

\section{Special Graph Algorithms}
In this section, we will discuss a special graph algorithm called Depth First Search (DFS). DFS is used to traverse or search a tree or graph data structure. It starts at a selected node and explores as far as possible along each branch before backtracking.
\FloatBarrier
\textbf{Algorithm Example: Depth First Search (DFS)}
\FloatBarrier
\begin{algorithm}
\caption{Depth First Search}
\begin{algorithmic}[1]
\Procedure{DFS}{node $v$}
\State Mark $v$ as visited
\For{each unvisited neighbor $u$ of $v$}
\State DFS($u$)
\EndFor
\EndProcedure
\end{algorithmic}
\end{algorithm}
\FloatBarrier
Let's consider a simple example of an undirected graph represented as an adjacency list:
\FloatBarrier
\begin{equation}
\begin{aligned}
\text{Graph Representation:} \\
\text{1: [2, 3]} \\
\text{2: [1, 4, 5]} \\
\text{3: [1]} \\
\text{4: [2]} \\
\text{5: [2]}
\end{aligned}
\end{equation}
\FloatBarrier
Applying DFS on this graph starting from node 1 would visit the nodes in the following order: 1, 2, 4, 5, 3.
\FloatBarrier
\textbf{Python Implementation}
\FloatBarrier
Here is the Python code implementation of DFS for the given graph:
\FloatBarrier
\begin{lstlisting}
def dfs(graph, node, visited):
    if node not in visited:
        visited.append(node)
        for neighbor in graph[node]:
            dfs(graph, neighbor, visited)
    return visited

graph = {
    1: [2, 3],
    2: [1, 4, 5],
    3: [1],
    4: [2],
    5: [2]
}
start_node = 1
visited_nodes = dfs(graph, start_node, [])
print("Visited nodes using DFS:", visited_nodes)
\end{lstlisting}
\FloatBarrier

\subsection{Topological Sort}
\subsubsection{Definition and Applications}
Topological sort is a linear ordering of vertices in a Directed Acyclic Graph (DAG) such that for every directed edge u -> v, vertex u comes before vertex v in the ordering. This ordering is useful in scheduling tasks or dependencies where one task must be completed before another.
\subsubsection{Implementing Topological Sort}
\FloatBarrier
\textbf{Algorithmic Example}
\FloatBarrier
Given a DAG represented as an adjacency list, we can perform a topological sort using Depth First Search (DFS). We start with an empty list for the topological ordering and a visited set. For each vertex, we recursively visit its neighbors and add the vertex to the ordering only after visiting all its neighbors.
\FloatBarrier
\begin{algorithm}
\caption{Topological Sort using DFS}
\begin{algorithmic}[1]
\State \textbf{function} topologicalSort(adjList, vertex, visited, ordering)
\State \hspace{\algorithmicindent} visited.add(vertex)
\For{neighbor \textbf{in} adjList[vertex]}
\If{neighbor \textbf{not in} visited}
\State topologicalSort(adjList, neighbor, visited, ordering)
\EndIf
\EndFor
\State ordering.insert(0, vertex)
\State \textbf{return}

\end{algorithmic}
\end{algorithm}
\FloatBarrier
\textbf{Python Code}
\FloatBarrier
\begin{lstlisting}
def topologicalSort(adjList):
    visited = set()
    ordering = []

    def dfs(vertex):
        visited.add(vertex)
        for neighbor in adjList[vertex]:
            if neighbor not in visited:
                dfs(neighbor)
        ordering.insert(0, vertex)

    for vertex in adjList:
        if vertex not in visited:
            dfs(vertex)

    return ordering
\end{lstlisting}
\FloatBarrier

\subsection{Algorithms for Finding Strong Components}

Strongly connected components in a directed graph are subsets of vertices where each vertex is reachable from every other vertex within the same subset. There are various algorithms to find strong components, with Kosaraju's algorithm being one of the most popular ones.

\subsubsection{Kosaraju's Algorithm}

Kosaraju's algorithm is based on the concept that if we reverse all the edges in a directed graph and perform a Depth First Search (DFS), the resulting forest of the DFS will consist of trees where each tree represents a strongly connected component.

Here is the detailed algorithm of Kosaraju's algorithm:
\FloatBarrier
\begin{algorithm}
\caption{Kosaraju's Algorithm}
\begin{algorithmic}

\State \textbf{Input:} Directed graph $G = (V, E)$
\State \textbf{Output:} List of strongly connected components

\State Perform a DFS on $G$ and store the finishing times of each vertex in a stack $S$
\State Reverse all the edges in $G$ to obtain $G_{rev}$
\State Initialize an empty list $SCC$
\State While $S$ is not empty:
\State \quad Pop a vertex $v$ from $S$
\State \quad If $v$ is not visited in $G_{rev}$:
\State \quad \quad Perform a DFS starting from $v$ in $G_{rev}$ to obtain a strongly connected component $C$
\State \quad \quad Add $C$ to $SCC$
\State Return $SCC$

\end{algorithmic}
\end{algorithm}
\FloatBarrier
\textbf{Algorithmic Example with Mathematical Detail}
\FloatBarrier
Let's consider the following directed graph $G = (V, E)$:

$V = \{A, B, C, D, E\}$

$E = \{(A, B), (B, D), (D, C), (C, A), (C, E), (E, D)\}$

Now, let's apply Kosaraju's algorithm to find the strongly connected components of $G$.
\FloatBarrier
\textbf{Python Code Equivalent}
\FloatBarrier
Here is the Python code equivalent to implement Kosaraju's algorithm:
\FloatBarrier
\begin{lstlisting}
# Python implementation of Kosaraju's Algorithm

def dfs(graph, vertex, visited, stack):
    visited.add(vertex)
    for neighbor in graph[vertex]:
        if neighbor not in visited:
            dfs(graph, neighbor, visited, stack)
    stack.append(vertex)

def transpose_graph(graph):
    transposed = {v: [] for v in graph.keys()}
    for vertex in graph:
        for neighbor in graph[vertex]:
            transposed[neighbor].append(vertex)
    return transposed

def kosaraju(graph):
    visited = set()
    stack = []

    for vertex in graph:
        if vertex not in visited:
            dfs(graph, vertex, visited, stack)

    transposed = transpose_graph(graph)
    visited.clear()
    scc = []

    while stack:
        vertex = stack.pop()
        if vertex not in visited:
            component = []
            dfs(transposed, vertex, visited, component)
            scc.append(component)

    return scc

# Example directed graph G
graph = {
    'A': ['B'],
    'B': ['D'],
    'C': ['A', 'E'],
    'D': ['C'],
    'E': ['D']
}

print("Strongly Connected Components:")
print(kosaraju(graph))
\end{lstlisting}
\FloatBarrier

\subsubsection{Tarjan's Strongly Connected Components Algorithm}
Tarjan's algorithm is used to find strongly connected components in a directed graph. Strongly connected components are subsets of vertices where each vertex is reachable from every other vertex within that subset.

\textbf{Algorithmic Example}
\FloatBarrier
\begin{algorithm}
\caption{Tarjan's Strongly Connected Components Algorithm}\label{tarjan}
\begin{algorithmic}[1]
\State Initialize $index=0$, $S=\emptyset$, $result=\emptyset$
\Function{StrongConnect}{v}
    \State $v.index \gets index$
    \State $v.lowlink \gets index$
    \State $index \gets index + 1$
    \State $S.push(v)$
    \For{each neighbour $w$ of $v$}
        \If{$w.index$ is not defined}
            \State \Call{StrongConnect}{$w$}
            \State $v.lowlink \gets \min(v.lowlink, w.lowlink)$
        \ElsIf{$w$ is in $S$}
            \State $v.lowlink \gets \min(v.lowlink, w.index)$
        \EndIf
    \EndFor
    \If{$v.lowlink = v.index$}
        \State Create a new component $C$
        \Repeat
            \State $w \gets S.pop()$
            \State Add $w$ to $C$
        \Until{$w = v$}
        \State Add $C$ to $result$
    \EndIf
\EndFunction
\end{algorithmic}
\end{algorithm}
\FloatBarrier
\textbf{Python Code}
\FloatBarrier
\begin{lstlisting}
# Python implementation of Tarjan's algorithm for SCC
def tarjan_SCC(graph):
    index = 0
    S = []
    result = []
    
    def strong_connect(v):
        nonlocal index
        v.index = index
        v.lowlink = index
        index += 1
        S.append(v)
        
        for w in v.neighbours:
            if not hasattr(w, 'index'):
                strong_connect(w)
                v.lowlink = min(v.lowlink, w.lowlink)
            elif w in S:
                v.lowlink = min(v.lowlink, w.index)
        
        if v.lowlink == v.index:
            C = []
            while True:
                w = S.pop()
                C.append(w)
                if w == v:
                    break
            result.append(C)
    
    for v in graph.vertices:
        if not hasattr(v, 'index'):
            strong_connect(v)
    
    return result
\end{lstlisting}
\FloatBarrier

\section{Minimum Cut and Graph Partitioning}

Minimum Cut in a graph refers to the partitioning of the vertices of a graph into two sets such that the number of edges between the two sets is minimized. This concept is essential in graph theory and has applications in various fields.

One of the popular algorithms used to find the minimum cut in a graph is the Karger's algorithm:
\FloatBarrier
\begin{algorithm}
\caption{Karger's Algorithm for Minimum Cut}
\begin{algorithmic}[1]
\Function{Karger}{$G$}
\While{$|V(G)| > 2$}
    \State Pick a random edge $(u, v)$ uniformly at random from $E(G)$
    \State Merge the vertices $u$ and $v$ into a single vertex
    \State Remove self-loops
\EndWhile
\State \Return the number of remaining edges
\EndFunction
\end{algorithmic}
\end{algorithm}
\FloatBarrier
To implement Karger's algorithm in Python, here is the corresponding code snippet:
\FloatBarrier
\begin{lstlisting}
import random

def karger(G):
    while len(G) > 2:
        u, v = random.choice(G.edges())  # Pick a random edge
        G.contract(u, v)  # Merge the vertices
        G.remove_self_loops()  # Remove self-loops
    return len(G.edges())

# Usage example
# G is a graph represented with vertices and edges
min_cut = karger(G)
print("Minimum Cut:", min_cut)
\end{lstlisting}
\FloatBarrier
This algorithm iteratively contracts edges in the graph until only two vertices remain, representing the two partitions. The number of remaining edges after the iterations gives the minimum cut value.

\subsection{The Concept of Minimum Cut}

The concept of minimum cut in a graph refers to the partitioning of the graph into two sets of vertices, such that the number of edges between the two sets is minimized. In other words, the minimum cut represents the smallest number of edges that need to be removed in order to disconnect the graph.

\textbf{Algorithmic Example}
\FloatBarrier
Here is an algorithmic example to find the minimum cut in a graph using the Ford-Fulkerson algorithm:
\FloatBarrier
\begin{algorithm}
\caption{Ford-Fulkerson Algorithm for Minimum Cut}
\begin{algorithmic}
\State Initialize residual graph $Gf$ with capacities and flows
\State Initialize empty cut set $S$
\While{there is a path $p$ from source $s$ to sink $t$ in $Gf$}
    \State Find the minimum residual capacity $c_f$ of path $p$
    \For{each edge $(u,v)$ in path $p$}
        \State Update the flow $f(u,v) = f(u,v) + c_f$ in $Gf$
        \State Update the reverse edge $(v,u)$ with $f(v,u) = f(v,u) - c_f$
        \If{vertex $v$ is reachable from source in $Gf$}
            \State Add vertex $v$ to set $S$
        \EndIf
    \EndFor
\EndWhile
\end{algorithmic}
\end{algorithm}
\FloatBarrier
\textbf{Python Implementation}
\FloatBarrier
Here is the Python code equivalent to the Ford-Fulkerson algorithm for finding the minimum cut in a graph:
\FloatBarrier
\begin{lstlisting}
def ford_fulkerson(G, s, t):
    # Initialize residual graph Gf with capacities and flows
    Gf = initialize_residual_graph(G)
    # Initialize empty cut set S
    S = set()
    # While there is a path p from source s to sink t in Gf
    while True:
        p = find_path(Gf, s, t)
        if not p:
            break
        # Find the minimum residual capacity cf of path p
        cf = min_residual_capacity(Gf, p)
        # Update the flow f(u,v) = f(u,v) + cf in Gf
        for u, v in zip(p[:-1], p[1:]):
            Gf[u][v]['flow'] += cf
            # Update the reverse edge (v,u) with f(v,u) = f(v,u) - cf
            Gf[v][u]['flow'] -= cf
            # If vertex v is reachable from source in Gf, add vertex v to set S
            if reachable(Gf, s, v):
                S.add(v)
    return S
\end{lstlisting}

\subsubsection{Definition and Importance}
Minimum Cut and Graph Partitioning are fundamental concepts in graph theory and computer science. A minimum cut in a graph is the smallest set of edges that, when removed, disconnects the graph into two or more components. It represents the minimal cost required to break the network into isolated components.

Graph partitioning involves dividing a graph into multiple subsets or partitions based on certain criteria. This partitioning is essential in various applications such as network optimization, clustering, and parallel computing. Finding optimal graph partitions can lead to efficient resource allocation and improved performance in various systems.
\FloatBarrier
\begin{algorithm}\label{algo:min_cut}
\caption{Minimum Cut Algorithm}
\begin{algorithmic}[1]
\State Let $G=(V,E)$ be a connected graph
\State Initialize $min\_cut = \infty$
\For{each vertex $v \in V$}
    \State Partition $G$ into two sets: $A = \{v\}$ and $B = V \setminus \{v\}$
    \State Compute the cut size between $A$ and $B$
    \State Update $min\_cut$ if the current cut size is smaller
\EndFor
\State \textbf{return} $min\_cut$
\end{algorithmic}
\end{algorithm}
\FloatBarrier
\begin{lstlisting}
    def min_cut(graph):
    min_cut = float('inf')
    for vertex in graph.vertices:
        A = {vertex}
        B = set(graph.vertices) - A
        cut_size = compute_cut_size(graph, A, B)
        min_cut = min(min_cut, cut_size)
    return min_cut
\end{lstlisting}
\subsubsection{Applications}
Minimum Cut and Graph Partitioning algorithms have numerous applications in various fields such as network design, image segmentation, clustering, and VLSI design. One common application is in community detection in social networks, where the goal is to partition the network into communities based on connectivity patterns.
\FloatBarrier
\textbf{Algorithmic Example: Karger's Min Cut Algorithm}
\FloatBarrier
Karger's Min Cut Algorithm is a randomized algorithm to find the minimum cut in an undirected graph.
\FloatBarrier
\begin{algorithm}
\caption{Karger's Min Cut Algorithm}
\begin{algorithmic}[1]
\Procedure{MinCut}{$G$}
\While{$|V(G)| > 2$}
\State Pick a random edge $(u, v)$ from $G$
\State Merge the nodes $u$ and $v$ into a single node $w$
\State Update the graph by removing self-loops and parallel edges
\EndWhile
\State \textbf{return} the number of remaining edges
\EndProcedure
\end{algorithmic}
\end{algorithm}
\FloatBarrier
The algorithm repeatedly contracts random edges until only two nodes remain in the graph, and the number of edges between them gives the minimum cut size.
\FloatBarrier
\textbf{Python Implementation of Karger's Min Cut Algorithm}
\FloatBarrier
Here is a Python implementation of Karger's Min Cut Algorithm:
\FloatBarrier
\begin{lstlisting}
import random

def min_cut(graph):
    while len(graph) > 2:
        # Randomly pick an edge (u, v) from the graph
        u = random.choice(list(graph.keys()))
        v = random.choice(graph[u])
        
        # Merge nodes u and v into a single node
        graph[u].extend(graph[v])
        
        for node in graph[v]:
            graph[node].remove(v)
            graph[node].append(u)
        
        while v in graph[u]:
            graph[u].remove(v)
        
        del graph[v]
    
    return len(list(graph.values())[0])
\end{lstlisting}
\FloatBarrier
This Python function takes a graph represented as an adjacency list and returns the size of the minimum cut.

\subsection{The Random Contraction Algorithm}
The Random Contraction Algorithm is a randomized algorithm used to find a minimum cut in a graph. The algorithm works by randomly "contracting" edges in the graph until only two vertices remain. The remaining edges between the two vertices form the minimum cut of the graph.
\FloatBarrier

\subsubsection{Algorithm Description}
\begin{algorithm}
\caption{Random Contraction Algorithm}
\begin{algorithmic}[1]
\Function{RandomContraction}{$G$}
    \While{$|V(G)| > 2$}
    \State Pick a random edge $(u,v)$ from $G$
    \State Merge vertices $u$ and $v$ into a single vertex
    \State Remove self-loops
    \EndWhile
\EndFunction
\end{algorithmic}
\end{algorithm}
\FloatBarrier
The algorithm operates by choosing a random edge at each iteration, merging its two endpoints into a single vertex, and removing self-loops until only two vertices remain.
\FloatBarrier
\textbf{Python code equivalent for the algorithmic example:}
\begin{lstlisting}

import random

def random_contraction(graph):
    while len(graph) > 2:
        # Pick a random edge from the graph
        u = random.choice(list(graph.keys()))
        v = random.choice(graph[u])

        # Merge vertices u and v into a single vertex
        graph[u].extend(graph[v])
        for node in graph[v]:
            graph[node].remove(v)
            graph[node].append(u)
        
        # Remove self-loops
        while u in graph[u]:
            graph[u].remove(u)
        
        del graph[v]

    return len(list(graph.values())[0])

# Example graph representation as an adjacency list
graph = {
    1: [2, 3, 4],
    2: [1, 3, 4],
    3: [1, 2, 4],
    4: [1, 2, 3]
}

print(random_contraction(graph))
    
\end{lstlisting}
\FloatBarrier
\subsubsection{Analysis and Applications}
The Random Contraction Algorithm is an efficient algorithm for finding the minimum cut in a graph. It works by repeatedly contracting random edges in the graph until there are only two vertices left, representing the two sides of the cut. This algorithm is simple to implement and has applications in various fields such as network analysis, image segmentation, and community detection.
\FloatBarrier
\begin{algorithm}
  \caption{Random Contraction Algorithm}
  \begin{algorithmic}[1]
    \State \textbf{Input:} Graph $G=(V,E)$
    \State \textbf{Output:} Minimum cut in $G$
    \While{$|V| > 2$}
      \State Choose a random edge $(u, v)$ from $E$
      \State Merge vertices $u$ and $v$ into a single vertex
      \State Remove self-loops
      \State Update edges by merging parallel edges
      \State Update $V$ and $E$
    \EndWhile
    \State \textbf{return} Remaining edges as the minimum cut
  \end{algorithmic}
\end{algorithm}

\FloatBarrier
\begin{lstlisting}
import random

def random_contraction_algorithm(graph):
    while len(graph) > 2:
        u, v = random.choice(graph.edges)
        graph.contract(u, v)
    return graph.edges

# Example graph
class Graph:
    def __init__(self):
        self.vertices = []
        self.edges = []

    def add_vertex(self, v):
        self.vertices.append(v)

    def add_edge(self, u, v):
        self.edges.append((u, v))

    def contract(self, u, v):
        for edge in self.edges:
            if edge[0] == v:
                self.edges.remove(edge)
                self.edges.append((u, edge[1]))
            elif edge[1] == v:
                self.edges.remove(edge)
                self.edges.append((u, edge[0]))
        self.vertices.remove(v)

# Instantiate graph and run algorithm
graph = Graph()
graph.add_vertex(1)
graph.add_vertex(2)
graph.add_vertex(3)
graph.add_vertex(4)
graph.add_edge(1, 2)
graph.add_edge(2, 3)
graph.add_edge(3, 4)
graph.add_edge(4, 1)

min_cut = random_contraction_algorithm(graph)
print("Minimum cut edges:", min_cut)
\end{lstlisting}
\FloatBarrier
\section{Randomized Algorithms in Graph Theory}
\subsection{Introduction to Randomized Algorithms}

Randomized algorithms in graph theory play a crucial role in solving various graph-related problems efficiently. These algorithms often use randomness in their decision-making process to achieve better performance or to simplify complex computations.
\FloatBarrier
\textbf{Example of a Randomized Algorithm}
One popular randomized algorithm in graph theory is Randomized Primality testing, also known as the Miller-Rabin primality test. The Miller-Rabin test is used to test whether a given number is prime or composite with high probability.
\subsection{Randomized Selection Algorithm}
\subsubsection{Algorithm Overview}
\FloatBarrier
\begin{algorithm}
\caption{Miller-Rabin Primality Test}
\begin{algorithmic}[1]
\Function{MillerRabin}{$n, k$}
    \State Let $d = n-1$
    \State Let $s = 0$
    \While{$d \mod 2 == 0$}
        \State $d \gets d / 2$
        \State $s \gets s + 1$
    \EndWhile
    \For{$i=1$ to $k$}
        \State Choose a random integer $a$ such that $2 \leq a \leq n-2$
        \State $x \gets a^d \mod n$
        \If{$x == 1$ or $x == n-1$}
            \State \textbf{continue}
        \EndIf
        \State $prime \gets$ \textbf{False}
        \For{$r=1$ to $s-1$}
            \State $x \gets x^2 \mod n$
            \If{$x == 1$}
                \State \textbf{return} \textbf{False} \Comment{$n$ is composite}
            \EndIf
            \If{$x == n-1$}
                \State \textbf{break}
            \EndIf
        \EndFor
        \If{$x \neq n-1$}
            \State \textbf{return} \textbf{False} \Comment{$n$ is composite}
        \EndIf
    \EndFor
    \State \textbf{return} \textbf{True} \Comment{$n$ is probably prime}
\EndFunction
\end{algorithmic}
\end{algorithm}
\FloatBarrier
The Python code for the Miller-Rabin Primality Test algorithm is as follows:
\FloatBarrier
\begin{lstlisting}
import random

def miller_rabin(n, k):
    d = n - 1
    s = 0
    while d % 2 == 0:
        d //= 2
        s += 1
    for _ in range(k):
        a = random.randint(2, n - 2)
        x = pow(a, d, n)
        if x == 1 or x == n - 1:
            continue
        prime = False
        for r in range(1, s):
            x = pow(x, 2, n)
            if x == 1:
                return False
            if x == n - 1:
                break
        if x != n - 1:
            return False
    return True
\end{lstlisting}
\FloatBarrier
\subsubsection{Analysis}
Randomized selection algorithms are fundamental in computer science and are often used to find the $k$th smallest element in an unsorted array. These algorithms are based on the principle of selecting a pivot element randomly and partitioning the array around this pivot.

The key idea behind randomized selection is that by choosing the pivot randomly, we can avoid the worst-case scenarios encountered in deterministic selection algorithms.

The analysis of randomized selection algorithms involves understanding the expected time complexity and the probability of selecting a good pivot. While the worst-case time complexity of randomized selection algorithms is linear, their expected time complexity is often sublinear.

Let $T(n)$ denote the expected time complexity of selecting the $k$th smallest element in an array of size $n$. The recurrence relation for the expected time complexity can be expressed as:

\[
T(n) = O(n) + T\left(\frac{n}{2}\right)
\]

This recurrence relation arises from the partitioning step of the algorithm, where we divide the array into two subarrays of approximately equal size.
\FloatBarrier
\textbf{Example}

Consider the following example of the randomized selection algorithm:
\FloatBarrier
\begin{lstlisting}[language=Python]
import random

def randomized_selection(arr, k):
    if len(arr) == 1:
        return arr[0]
    
    pivot = random.choice(arr)
    smaller = [x for x in arr if x < pivot]
    equal = [x for x in arr if x == pivot]
    larger = [x for x in arr if x > pivot]
    
    if k < len(smaller):
        return randomized_selection(smaller, k)
    elif k < len(smaller) + len(equal):
        return equal[0]
    else:
        return randomized_selection(larger, k - len(smaller) - len(equal))
\end{lstlisting}
\FloatBarrier

This Python code demonstrates the implementation of the randomized selection algorithm. By selecting the pivot randomly and partitioning the array into smaller, equal, and larger subarrays, the algorithm efficiently finds the $k$th smallest element in the array.

\subsubsection{Applications in Graph Algorithms}

Randomized selection algorithms are powerful tools in graph theory, known for their efficiency and effectiveness in solving complex problems. Here's how they are commonly applied:

\textbf{Randomized Minimum Cut Algorithm}
This approach uses a randomized algorithm to determine a graph's minimum cut efficiently. By randomly selecting edges and contracting them until only two vertices are left, the algorithm finds the minimum cut as the edges connecting these final two vertices.

\textbf{Randomized Approximation Algorithms}
For problems like the maximum cut, graph coloring, and traveling salesman problem, randomized algorithms offer approximation solutions that are often near-optimal. They are particularly valuable for handling large-scale graphs where exact solutions would be computationally prohibitive.

\textbf{Randomized Graph Partitioning}
In graph partitioning, the objective is to split a graph into subgraphs while minimizing the edges cut. Randomized selection algorithms facilitate this by randomly picking vertices or edges to form partitions, making the process efficient and scalable.

\textbf{Randomized Sampling}
These algorithms are also utilized for sampling within large graphs. Randomly selected vertices or edges can provide insights into the overall structure of the graph, aiding tasks like community detection, anomaly identification, and visualization.

These applications underscore the broad utility of randomized selection algorithms in graph theory, offering solutions that balance efficiency with accuracy across various graph-related challenges.


\section{Graph Search Strategies}
Graph search strategies are algorithms used to traverse or search a graph in order to find a particular vertex or a path between two vertices. Some common graph search strategies include Depth-First Search (DFS) and Breadth-First Search (BFS).

\subsection{Advanced BFS and DFS Techniques}
\textbf{Breadth-First Search:}
\FloatBarrier
In this section, we will explore advanced techniques related to Breadth First Search (BFS).
\FloatBarrier
\textbf{Algorithmic Example}
\FloatBarrier
\begin{algorithm}
\caption{Advanced BFS Algorithm}
\begin{algorithmic}[1]
\State Initialize Queue $Q$ with start node
\State Mark start node as visited
\While{Q is not empty}
    \State Dequeue current node $u$ from $Q$
    \For{each neighbor $v$ of $u$}
        \If{$v$ is not visited}
            \State Mark $v$ as visited
            \State Enqueue $v$ into $Q$
        \EndIf
    \EndFor
\EndWhile
\end{algorithmic}
\end{algorithm}
\FloatBarrier
The above algorithm demonstrates how to perform an advanced BFS traversal on a graph.
\FloatBarrier
\textbf{Python Implementation}
\FloatBarrier
\begin{lstlisting}
# Python Implementation of Advanced BFS Algorithm
from collections import deque

def advanced_bfs(graph, start):
    queue = deque([start])
    visited = set([start])

    while queue:
        current_node = queue.popleft()
        for neighbor in graph[current_node]:
            if neighbor not in visited:
                visited.add(neighbor)
                queue.append(neighbor)
\end{lstlisting}
\FloatBarrier
\textbf{Depth First Search}
\FloatBarrier
Depth First Search (DFS) is a fundamental graph traversal algorithm that explores as far as possible along each branch before backtracking. In this section, we delve into some advanced techniques to enhance the efficiency and utility of DFS.
\FloatBarrier
\textbf{Algorithmic Example: Depth First Search with Backtracking}
\FloatBarrier
\begin{algorithm}
    \caption{DFS with Backtracking}
    \begin{algorithmic}[1]
        \Function{DFS}{node}
            \State Mark $node$ as visited
            \For{each neighbor $neigh$ of $node$}
                \If{$neigh$ is not visited}
                    \State Recursive call: \textsc{DFS}($neigh$)
                \EndIf
            \EndFor
        \EndFunction
    \end{algorithmic}
\end{algorithm}
\FloatBarrier
To demonstrate the above algorithm, consider a simple graph represented by an adjacency list:
\[
\{A: [B, C], B: [D], C: [E], D: [F], E: [F], F: []\}
\]
\FloatBarrier
The Python equivalent code for the above algorithm is as follows:
\FloatBarrier
\begin{lstlisting}
def dfs(node, visited):
    visited[node] = True
    for neigh in graph[node]:
        if not visited[neigh]:
            dfs(neigh, visited)

# Example usage with the given graph
graph = {'A': ['B', 'C'], 'B': ['D'], 'C': ['E'], 'D': ['F'], 'E': ['F'], 'F': []}
visited = {node: False for node in graph}
dfs('A', visited)
\end{lstlisting}
\FloatBarrier
\subsection{Heuristic Search Algorithms in Graphs}
Heuristic search algorithms are a class of algorithms used to solve problems in a more efficient manner by using heuristics. In graph theory, heuristic search algorithms play a crucial role in finding optimal or near-optimal solutions for various graph-related problems. These algorithms leverage heuristic information to guide the search process towards solutions.
\subsubsection{A* Search Algorithm}
One commonly used heuristic search algorithm in graph theory is the A* algorithm. A* is an informed search algorithm that uses both the cost to reach a node (denoted by g) and a heuristic estimation of the cost to reach the goal node from the current node (denoted by h) to determine the next node to explore.

\FloatBarrier
\textbf{Algorithmic example with mathematical detail:}
\FloatBarrier
\begin{algorithm}
    \caption{A* Algorithm}
    \begin{algorithmic}[1]
        \State Initialize an open list with the initial node
        \State Initialize an empty set of explored nodes
        \While{open list is not empty}
            \State Choose the node with the lowest value of f = g + h from the open list
            \State Move this node to the explored nodes set
            \If{current node is the goal node}
                \State \textbf{return} path
            \EndIf
            \For{each neighbor of the current node}
                \If{neighbor is not in the open list or explored nodes}
                    \State Calculate g and h values for the neighbor
                    \State Add the neighbor to the open list
                \EndIf
            \EndFor
        \EndWhile
    \end{algorithmic}
\end{algorithm}
\FloatBarrier
\textbf{Python code equivalent for the A* algorithm:}
\FloatBarrier
\begin{lstlisting}

def astar_algorithm(graph, start, goal):
    open_list = {start}
    explored = set()
    g = {node: float('inf') for node in graph}
    g[start] = 0
    h = {node: heuristic(node, goal) for node in graph}

    while open_list:
        current = min(open_list, key=lambda node: g[node] + h[node])
        open_list.remove(current)
        explored.add(current)

        if current == goal:
            return reconstruct_path(goal)

        for neighbor in graph[current]:
            if neighbor not in explored:
                tentative_g = g[current] + graph[current][neighbor]
                if tentative_g < g.get(neighbor, float('inf')):
                    g[neighbor] = tentative_g
                    open_list.add(neighbor)

    return None
    \end{lstlisting}
\FloatBarrier

\subsubsection{Greedy Best-First Search}
Greedy Best-First Search is a variation of Best-First Search algorithm that prioritizes nodes based on a heuristic evaluation function. At each step, it selects the node that appears to be the most promising based on the heuristic.
\FloatBarrier
\textbf{Algorithm}
\FloatBarrier
\begin{algorithm}
\caption{Greedy Best-First Search}
\begin{algorithmic}
\State \textbf{Input:} Graph $G$, start node $s$, heuristic function $h$
\State \textbf{Output:} Path from $s$ to goal node

\State Initialize priority queue $Q$ with $s$
\State Initialize an empty set $visited$
\While{$Q$ is not empty}
    \State $current \gets Q$.pop()
    \If{$current$ is goal node}
        \State \textbf{return} path to $current$
    \EndIf
    \State Add $current$ to $visited$
    \For{each neighbor $n$ of $current$}
        \If{$n$ is not in $visited$}
            \State Calculate priority $p$ using heuristic function: $p = h(n)$
            \State Add $n$ to $Q$ with priority $p$
        \EndIf
    \EndFor
\EndWhile
\State \textbf{return} No path found
\end{algorithmic}
\end{algorithm}
\FloatBarrier
\textbf{Python Code}
\FloatBarrier
Here is the python code equivalent for the Greedy Best-First Search algorithm:
\FloatBarrier
\begin{lstlisting}
from queue import PriorityQueue

def greedy_best_first_search(graph, start, goal, heuristic):
    pq = PriorityQueue()
    pq.put((0, start))

    visited = set()

    while not pq.empty():
        current_cost, current_node = pq.get()

        if current_node == goal:
            return current_cost

        visited.add(current_node)

        for neighbor in graph[current_node]:
            if neighbor not in visited:
                priority = heuristic(neighbor)
                pq.put((priority, neighbor))

    return None
\end{lstlisting}
\FloatBarrier

\section{Graph Path Finding Algorithms}
Graph path finding algorithms are a set of techniques to find the shortest path between nodes in a graph. These algorithms are essential in various applications like navigation systems, network routing, and social network analysis.

One of the well-known algorithms for solving shortest path problems is Dijkstra's algorithm. This algorithm finds the shortest path from a starting node to all other nodes in a weighted graph. It uses a priority queue to greedily select the node with the smallest distance at each step until all nodes have been visited.

\subsection{Shortest Path Algorithms}
Shortest path algorithms are crucial in graph theory for finding the quickest route between two vertices in a weighted graph. These algorithms are widely used in network routing, transportation planning, and other optimization scenarios where the goal is to minimize travel cost or distance.

\textbf{Dijkstra's Algorithm}
One of the most renowned shortest path algorithms is Dijkstra's algorithm. Developed by Edsger W. Dijkstra in 1956, it efficiently computes the shortest paths from a single source vertex to all other vertices in a graph with non-negative edge weights. Dijkstra's algorithm uses a priority queue to keep track of vertex distances and iteratively updates the paths and distances until all vertices are processed.

\textbf{Bellman-Ford Algorithm}
For graphs that include negative edge weights, the Bellman-Ford algorithm is more suitable. It iterates over the graph's edges to minimize the path distances, also providing the capability to detect negative cycles, which are crucial for understanding feasibility and stability in networks.

\textbf{Specialized Algorithms}
Other algorithms like the A* algorithm cater to graphs with non-negative integer weights by utilizing heuristics to speed up the search, whereas the Floyd-Warshall algorithm addresses scenarios with negative weights by finding shortest paths between all pairs of vertices.

\subsubsection{Dijkstra's Algorithm in Detail}
Dijkstra's algorithm operates as follows:
- Initialize distances: Set the distance from the source to itself to zero and all others to infinity.
- Use a priority queue: Store vertices by distance from the source.
- Relax edges: For the vertex with the shortest distance in the queue, update the distances for its adjacent vertices.
- Repeat: Continue until all vertices are processed.

The result is a set of shortest paths from the source to all vertices, enabling efficient route planning and network design.

Overall, shortest path algorithms like Dijkstra's provide essential tools for designing and optimizing systems across various fields, enhancing the efficiency of resources and operational planning.


Here is the implementation of Dijkstra's algorithm in Python:
\FloatBarrier
\begin{lstlisting}
    import heapq

def dijkstra(graph, source):
    # Initialize distances to all vertices as infinity
    distances = {vertex: float('infinity') for vertex in graph}
    distances[source] = 0
    
    # Priority queue to store vertices sorted by distance
    pq = [(0, source)]
    
    while pq:
        current_distance, current_vertex = heapq.heappop(pq)
        
        # Skip if the current distance is greater than the known distance
        if current_distance > distances[current_vertex]:
            continue
            
        # Iterate over neighbors of the current vertex
        for neighbor, weight in graph[current_vertex].items():
            distance = current_distance + weight
            
            # Relax the edge if a shorter path is found
            if distance < distances[neighbor]:
                distances[neighbor] = distance
                heapq.heappush(pq, (distance, neighbor))
    
    return distances
\end{lstlisting}
In this implementation, graph is a dictionary representing the weighted graph, where keys are vertices and values are dictionaries of neighbors and their corresponding edge weights. source is the starting vertex for which the shortest paths are to be computed. The function returns a dictionary containing the shortest distances from the source vertex to all other vertices in the graph.

\subsubsection{Bellman-Ford Algorithm}
The Bellman-Ford algorithm, developed by Richard Bellman and Lester Ford Jr. in the 1950s, finds shortest paths from a single source vertex to all other vertices in a weighted graph. This method is particularly valuable because it can handle graphs with negative weight edges, a feature not supported by Dijkstra's algorithm.

\textbf{How It Works:}
The Bellman-Ford algorithm operates through a process called edge relaxation, which iteratively updates the shortest path estimates:

- \textbf{Initialization:} Set the distance from the source vertex to itself to zero and all other vertex distances to infinity.

- \textbf{Edge Relaxation:} For \(|V|-1\) iterations—where \(|V|\) is the number of vertices—relax all the edges. This involves checking each edge and, if a shorter path is found, updating the distance estimate for the vertex at the end of that edge.

- \textbf{Negative Cycle Detection:} After \(|V|-1\) cycles, any further ability to relax an edge indicates a negative weight cycle in the graph.

\textbf{Practical Implications:}
Though less efficient than Dijkstra's algorithm due to its need to relax edges multiple times, Bellman-Ford is crucial for applications where negative weights are involved and there's a need to ensure no negative cycles exist. The algorithm guarantees that the shortest paths calculated after \(|V|-1\) iterations are optimal unless a negative cycle affects the distances.

By providing a method to handle complex edge weight configurations, the Bellman-Ford algorithm remains a fundamental tool for network design and analysis, ensuring robust route planning even under challenging conditions.


Here's the Python implementation of the Bellman-Ford algorithm:
\FloatBarrier
\begin{lstlisting}
    def bellman_ford(graph, source):
    # Step 1: Initialize distances
    distances = {vertex: float('infinity') for vertex in graph}
    distances[source] = 0
    
    # Step 2: Relax edges for |V| - 1 iterations
    for _ in range(len(graph) - 1):
        for u, edges in graph.items():
            for v, weight in edges.items():
                if distances[u] + weight < distances[v]:
                    distances[v] = distances[u] + weight
    
    # Step 3: Check for negative weight cycles
    for u, edges in graph.items():
        for v, weight in edges.items():
            if distances[u] + weight < distances[v]:
                raise ValueError("Graph contains negative weight cycle")
    
    return distances
\end{lstlisting}
\FloatBarrier
In this implementation, graph is a dictionary representing the weighted graph, where keys are vertices and values are dictionaries of neighbors and their corresponding edge weights. source is the starting vertex for which the shortest paths are to be computed. The function returns a dictionary containing the shortest distances from the source vertex to all other vertices in the graph, or raises an error if a negative weight cycle is detected.

\subsection{All-Pairs Shortest Paths}
The "All-Pairs Shortest Paths" problem is pivotal in graph theory and deals with finding the shortest path between every pair of vertices in a weighted graph. This task is essential in various applications like network routing, transportation planning, and social network analysis, where understanding the shortest distances between all nodes is crucial for optimization and planning.

\textbf{Floyd-Warshall Algorithm}
A key tool for solving this problem is the Floyd-Warshall algorithm, renowned for its efficiency with dense graphs. It utilizes dynamic programming to iteratively refine the shortest path distances between all vertex pairs. The algorithm keeps a matrix that records these distances and updates it by considering whether taking paths through intermediate vertices offers a shorter path between any two vertices.

\textbf{Johnson's Algorithm}
For sparser graphs, Johnson's algorithm provides a more suitable alternative. It cleverly merges the approaches of Dijkstra’s and Bellman-Ford algorithms, accommodating graphs with negative edge weights while efficiently handling sparse connections.

\textbf{Applications and Practical Importance}
From optimizing data flow between routers in communication networks to reducing travel times in urban planning, the solutions to the All-Pairs Shortest Paths problem play a foundational role. They not only enhance efficiency in real-world networks but also contribute to more robust and effective system designs.

\subsubsection{Floyd-Warshall Algorithm in Detail}
Independently developed by Bernard Roy and Stephen Warshall in the early 1960s and refined by Robert Floyd, the Floyd-Warshall algorithm is a cornerstone in computing shortest paths in weighted graphs. This algorithm performs particularly well in graphs that can have negative edge weights but must be free of negative weight cycles.

\textbf{Operation:}
The algorithm updates a distance matrix by systematically checking all possible paths through each vertex to see if a shorter path exists, using a methodical dynamic programming approach. It requires \(O(V^3)\) time complexity, where \(V\) is the number of vertices, making it especially suitable for graphs where the vertex count isn't prohibitively large.

\textbf{Significance:}
This method’s ability to compute shortest paths between all pairs of vertices simultaneously makes it invaluable for analyses where every inter-node distance is critical, supporting thorough and efficient evaluations of complex networks.


Here are the main steps of the Floyd-Warshall algorithm:

\begin{enumerate}
    \item Initialize a $|V| \times |V|$ matrix distances to represent the shortest distances between all pairs of vertices. Initially, the matrix is filled with the weights of the edges in the graph, and the diagonal elements are set to 0.
    \item For each intermediate vertex $k$ from 1 to $|V|$, iterate over all pairs of vertices $(i, j)$ and update the \texttt{distances[i][j]} entry if the path through vertex $k$ produces a shorter distance.
    \item After the iterations are complete, the \texttt{distances} matrix will contain the shortest distances between all pairs of vertices in the graph.
\end{enumerate}

\textbf{Python Implementation:}
\FloatBarrier
\begin{lstlisting}
    def floyd_warshall(graph):
    # Initialize the distance matrix
    distances = {u: {v: float('inf') for v in graph} for u in graph}
    for u in graph:
        distances[u][u] = 0
    for u in graph:
        for v, weight in graph[u].items():
            distances[u][v] = weight
    
    # Update distances using Floyd-Warshall algorithm
    for k in graph:
        for u in graph:
            for v in graph:
                distances[u][v] = min(distances[u][v], distances[u][k] + distances[k][v])
    
    return distances
\end{lstlisting}
In this implementation, graph is a dictionary representing the weighted graph, where keys are vertices and values are dictionaries of neighbors and their corresponding edge weights. The function returns a nested dictionary containing the shortest distances between all pairs of vertices in the graph.

\subsubsection{Johnson's Algorithm}
Johnson's algorithm is used to find the shortest paths between all pairs of vertices in a weighted graph, even in the presence of negative edge weights and negative weight cycles. It was developed by Donald B. Johnson in 1977. Johnson's algorithm combines Dijkstra's algorithm with the Bellman-Ford algorithm to handle negative edge weights.

Here's how Johnson's algorithm works:

\begin{enumerate}
    \item Add a new vertex, called the "dummy vertex," to the graph and connect it to all other vertices with zero-weight edges. This transforms the original graph into a new graph with no negative edge weights.
    \item Run the Bellman-Ford algorithm on the new graph with the dummy vertex as the source. This step detects negative weight cycles, if any.
    \item If the Bellman-Ford algorithm detects a negative weight cycle, terminate the algorithm and report that the graph contains a negative weight cycle.
    \item If there are no negative weight cycles, reweight the edges of the original graph using vertex potentials computed during the Bellman-Ford algorithm.
    \item For each vertex in the original graph, run Dijkstra's algorithm to find the shortest paths to all other vertices. Use the reweighted edge weights obtained in step 4.
    \item After running Dijkstra's algorithm for all vertices, compute the shortest paths between all pairs of vertices based on the distances obtained.
\end{enumerate}
\FloatBarrier
\textbf{Python Implementation:}
\FloatBarrier
\begin{lstlisting}
    import heapq
from collections import defaultdict

def bellman_ford(graph, source):
    distance = {v: float('inf') for v in graph}
    distance[source] = 0
    for _ in range(len(graph) - 1):
        for u, neighbors in graph.items():
            for v, weight in neighbors.items():
                distance[v] = min(distance[v], distance[u] + weight)
    return distance

def dijkstra(graph, source):
    distances = {v: float('inf') for v in graph}
    distances[source] = 0
    heap = [(0, source)]
    while heap:
        dist_u, u = heapq.heappop(heap)
        if dist_u > distances[u]:
            continue
        for v, weight in graph[u].items():
            if distances[u] + weight < distances[v]:
                distances[v] = distances[u] + weight
                heapq.heappush(heap, (distances[v], v))
    return distances

def johnson(graph):
    # Step 1: Add a dummy vertex and connect it to all other vertices with zero-weight edges
    dummy = 'dummy'
    graph[dummy] = {v: 0 for v in graph}

    # Step 2: Run Bellman-Ford algorithm to detect negative weight cycles and compute vertex potentials
    potentials = bellman_ford(graph, dummy)
    if any(potentials[v] < 0 for v in graph):
        return "Negative weight cycle detected"

    # Step 3: Reweight the edges using vertex potentials
    for u, neighbors in graph.items():
        for v in neighbors:
            graph[u][v] += potentials[u] - potentials[v]

    # Step 4: Run Dijkstra's algorithm for each vertex to find shortest paths to all other vertices
    shortest_paths = {}
    for u in graph:
        shortest_paths[u] = dijkstra(graph, u)

    # Step 5: Restore original distances using vertex potentials
    for u, distances in shortest_paths.items():
        for v in distances:
            shortest_paths[u][v] += potentials[v] - potentials[u]

    return shortest_paths
\end{lstlisting}
\FloatBarrier
This implementation returns a dictionary containing the shortest paths between all pairs of vertices in the graph. If a negative weight cycle is detected, the function returns the message "Negative weight cycle detected."

\section{Network Flow and Matching}
Network Flow and Matching are fundamental concepts in graph theory and algorithms. In the context of network flow, the goal is to determine how to optimally transport items through a network, while in matching, the objective is to find pairings or matchings between elements of different sets that satisfy certain criteria.

\subsection{Maximum Flow Problem}
The Maximum Flow Problem is a fundamental problem in graph theory and network flow optimization. Given a directed graph with capacities on the edges, the goal is to find the maximum amount of flow that can be sent from a source node to a sink node. This problem has various applications in transportation, communication networks, and more.

\textbf{Algorithmic Example}
\FloatBarrier
The Ford-Fulkerson algorithm is a popular algorithm for solving the Maximum Flow Problem. It iteratively finds augmenting paths from the source to the sink and increases the flow along these paths until no more augmenting paths can be found. The algorithm terminates when no more augmenting paths exist, and the flow is maximized.
\FloatBarrier
\begin{algorithm}
\caption{Ford-Fulkerson Algorithm}
\begin{algorithmic}
\Require Directed graph $G$, source node $s$, sink node $t$
\Ensure Maximum flow value $max\_flow$
\Function{FordFulkerson}{$G, s, t$}
    \State Initialize flow $f$ on all edges to 0
    \State $max\_flow \gets 0$
    \While{there exists an augmenting path $p$ from $s$ to $t$ in residual graph $G_f$}
        \State Find the bottleneck capacity $c_f(p)$ of path $p$
        \State Augment flow $f$ along path $p$
        \State Update residual capacities in $G_f$
        \State Update $max\_flow$ with $c_f(p)$
    \EndWhile
    \State \Return $max\_flow$
\EndFunction
\end{algorithmic}
\end{algorithm}
\FloatBarrier
\textbf{Python Implementation:}
\FloatBarrier
\begin{lstlisting}[language=Python, caption=Ford-Fulkerson Algorithm in Python]
def ford_fulkerson(G, s, t):
    def dfs(node, min_capacity):
        if node == t:
            return min_capacity
        for edge, residual in G[node]:
            if residual > 0:
                flow = dfs(edge, min(min_capacity, residual))
                if flow > 0:
                    G[node].append((edge, flow))
                    return flow
        return 0

    max_flow = 0
    while True:
        flow = dfs(s, float('inf'))
        if flow == 0:
            break
        max_flow += flow

    return max_flow
\end{lstlisting}
\subsubsection{Ford-Fulkerson Method}
The Ford-Fulkerson method is a technique for computing the maximum flow in a flow network. It operates by incrementing the flow along augmenting paths from the source to the sink until no more augmenting paths exist.
\FloatBarrier
\textbf{Algorithmic Example}
\FloatBarrier
Let $G = (V, E)$ be a flow network with capacities $c(u, v)$ where $(u, v) \in E$. The Ford-Fulkerson algorithm finds the maximum flow from a source node $s$ to a sink node $t$ in $G$.
\FloatBarrier
\begin{algorithm}
\caption{Ford-Fulkerson Algorithm}
\begin{algorithmic}
\State Initialize the flow on each edge to 0
\While{there exists an augmenting path from $s$ to $t$}
    \State Find an augmenting path $P$ in the residual graph $G_f$
    \State Compute the residual capacity $c_f(u, v)$ for each edge $(u, v) \in P$
    \State Update the flow along the path $P$
\EndWhile
\end{algorithmic}
\end{algorithm}
\FloatBarrier
\textbf{Python Code}
\FloatBarrier
Below is a Python implementation of the Ford-Fulkerson algorithm:
\FloatBarrier
\begin{lstlisting}
def ford_fulkerson(graph, source, sink):
    def dfs(graph, s, t, path, visited):
        if s == t:
            return path
        for node, capacity in graph[s].items():
            if node not in visited and capacity > 0:
                visited.add(node)
                result = dfs(graph, node, t, path + [(s, node)], visited)
                if result is not None:
                    return result
        return None
    
    max_flow = 0
    path = dfs(graph, source, sink, [], set())
    
    while path is not None:
        flow = min(graph[u][v] for u, v in path)
        max_flow += flow
        for u, v in path:
            graph[u][v] -= flow
            graph[v][u] += flow
        path = dfs(graph, source, sink, [], set())
    
    return max_flow
\end{lstlisting}

\subsubsection{Edmonds-Karp Algorithm}
The Edmonds-Karp Algorithm is a specific implementation of the Ford-Fulkerson method for computing the maximum flow in a flow network. It uses the concept of augmenting paths to iteratively find the maximum flow from a source vertex to a sink vertex in a flow network.
\FloatBarrier
\textbf{Algorithmic example for the Edmonds-Karp Algorithm:}
\FloatBarrier
\begin{algorithm}
\caption{Edmonds-Karp Algorithm}
\begin{algorithmic}[1]
\State Initialize flow network $G$ with capacities and flow values
\State Initialize maximum flow $0$
\While{there exists an augmenting path $p$ from source $s$ to sink $t$ in residual network $G_f$}
    \State Find the minimum residual capacity along path $p$ as $c_f(p)$
    \State Update the flow along path $p$ by adding $c_f(p)$ to $f$
    \State Update the residual capacities and flow values of edges in $p$
    \State Update the maximum flow with $c_f(p)$
\EndWhile
\end{algorithmic}
\end{algorithm}
\FloatBarrier
\textbf{Python code equivalent for the Edmonds-Karp Algorithm:}
\FloatBarrier
\begin{lstlisting}
def edmonds_karp(G, source, sink):
    flow = 0
    residual_graph = compute_residual_graph(G)
    while True:
        path = bfs(residual_graph, source, sink)
        if not path:
            break
        flow_augment = min(residual_graph[u][v]['capacity'] for u, v in path)
        for u, v in path:
            residual_graph[u][v]['capacity'] -= flow_augment
            residual_graph[v][u]['capacity'] += flow_augment
        flow += flow_augment
    return flow
        
\end{lstlisting}
\subsection{Minimum Cost Flow Problem}
The minimum cost flow problem is a network optimization problem where the goal is to find the cheapest way to send a certain amount of flow through a flow network. Each edge in the network has a capacity (maximum flow that can pass through it) and a cost per unit of flow. The objective is to minimize the total cost of sending the required flow from a source node to a sink node.
\FloatBarrier
\textbf{Algorithmic Example}
\FloatBarrier
Let's consider a simple example where we have a flow network represented by a directed graph with nodes $s$ as the source and $t$ as the sink. Each edge in the graph has a capacity and a cost associated with it. We want to find the minimum cost to send a certain amount of flow $F$ from $s$ to $t$.

We can solve this problem using the Minimum Cost Flow algorithm, which is based on the Ford-Fulkerson method with the Bellman-Ford shortest path algorithm for finding the minimum cost.
\FloatBarrier
\begin{algorithm}
\caption{Minimum Cost Flow Algorithm}
\begin{algorithmic}[1]
\State Initialize flow on each edge to 0
\State While there exists a path from $s$ to $t$ with available capacity:
\State \hspace{10mm} Find the path with the minimum cost using Bellman-Ford
\State \hspace{10mm} Update the flow along the path
\State Calculate the total cost as the sum of costs of all edges with non-zero flow
\end{algorithmic}
\end{algorithm}
\FloatBarrier
\textbf{Python Code}
\FloatBarrier
Here is the Python implementation of the Minimum Cost Flow algorithm for the described example:
\FloatBarrier
\begin{lstlisting}
import networkx as nx

def min_cost_flow(graph, source, sink, flow):
    flow_cost, flow_dict = nx.network_simplex(graph, demand={source: -flow, sink: flow})
    return flow_cost

G = nx.DiGraph()
G.add_edge('s', 'a', capacity=10, weight=1)
G.add_edge('s', 'b', capacity=5, weight=2)
G.add_edge('a', 'b', capacity=3, weight=3)
G.add_edge('a', 't', capacity=8, weight=4)
G.add_edge('b', 't', capacity=7, weight=5)

min_cost = min_cost_flow(G, 's', 't', 10)
print(f"Minimum cost for flow of 10 units from 's' to 't': {min_cost}")
    
\end{lstlisting}
\subsubsection{Applications and Solutions}
The minimum cost flow problem is a fundamental problem in optimization theory with various practical applications. Some common applications include:
\begin{itemize}
    \item Transportation and logistics planning: Optimizing the flow of goods through a network while minimizing transportation costs.
    \item Communication network design: Finding the most cost-effective way to route data through a network.
    \item Supply chain management: Determining the optimal distribution of resources from suppliers to consumers.
\end{itemize}
\FloatBarrier

To solve the minimum cost flow problem, algorithms such as the Network Simplex Method or the Successive Shortest Path Algorithm are commonly used. These algorithms efficiently find the flow that minimizes the total cost while satisfying the capacity constraints of the network.

\textbf{Algorithmic Example with Mathematical Detail:}
\FloatBarrier
\begin{algorithm}
\caption{Minimum Cost Flow Algorithm}
\begin{algorithmic}[1]
\State Initialize flow values on edges
\State Repeat until convergence:
\State \quad Solve the shortest path problem to find a negative cost cycle
\State \quad Update flow along the cycle
\end{algorithmic}
\end{algorithm}
\FloatBarrier
\textbf{Python Code Equivalent for the Algorithmic Example:}
\FloatBarrier
\begin{lstlisting}
def minimum_cost_flow_algorithm():
    # Initialize flow values on edges
    # Repeat until convergence
    while not converged:
        # Solve the shortest path problem to find a negative cost cycle
        negative_cycle = find_negative_cycle()
        # Update flow along the cycle
        update_flow(negative_cycle)
    
\end{lstlisting}
\section{Conclusion}
\subsection{Summary of Graph Algorithms}
Graph algorithms are used to analyze relationships between elements in a graph. They can be used to solve various problems such as finding the shortest path, determining connectivity, identifying cycles, and more. Here, we provide an example of a basic graph algorithm.

\subsection{Challenges in Graph Algorithm Research}
Graph algorithm research faces several challenges, including:

1. \textbf{Scalability}: As graphs grow larger in scale and complexity, algorithm efficiency becomes crucial. Developing algorithms that can handle massive graphs efficiently is a significant challenge.

2. \textbf{Dynamic Graphs}: Real-world graphs often evolve over time, with edges and vertices being added or removed dynamically. Designing algorithms that can adapt to such changes and maintain correctness is a challenging area of research.

3. \textbf{Complexity Analysis}: Analyzing the complexity of graph algorithms and determining their time and space complexity is not always straightforward, especially for algorithms dealing with highly connected or dense graphs.

\subsection{Future Directions in Graph Theory and Algorithm Research}
Graph theory and algorithms have been extensively studied and applied in various fields like computer science, mathematics, and network analysis. There are several exciting future directions for research in this area. One such direction is the development of algorithms for handling massive graphs efficiently. With the growth of data, analyzing and processing large-scale graphs poses a significant challenge.

Another promising avenue is the exploration of new graph properties and structures that can lead to the development of innovative graph algorithms. Additionally, integrating graph theory with other fields such as machine learning and data science opens up new possibilities for advanced applications and research.


\section{Further Reading and Resources}
To deepen your understanding of graph algorithms, there are numerous resources available that range from academic papers and books to online tutorials and open-source libraries. This section provides a comprehensive guide to these resources, helping you explore the topic in greater detail and apply what you have learned to real-world problems.

\subsection{Key Papers and Books on Graph Algorithms}
A solid foundation in graph algorithms can be built by studying some of the key papers and books in the field. These resources provide both theoretical insights and practical applications.

\begin{itemize}
    \item \textbf{“Graph Theory” by Reinhard Diestel} - This book is a comprehensive introduction to graph theory, covering fundamental concepts and theorems.
    \item \textbf{“Introduction to Algorithms” by Thomas H. Cormen, Charles E. Leiserson, Ronald L. Rivest, and Clifford Stein} - Often referred to as CLRS, this book is a must-read for anyone studying algorithms. It includes detailed sections on graph algorithms, such as breadth-first search, depth-first search, and shortest paths.
    \item \textbf{“Network Flows: Theory, Algorithms, and Applications” by Ravindra K. Ahuja, Thomas L. Magnanti, and James B. Orlin} - This book delves into the theory and algorithms related to network flows, providing both theoretical and practical perspectives.
    \item \textbf{“The Annotated Turing” by Charles Petzold} - While not exclusively about graph algorithms, this book offers an insightful look into the foundations of computer science, which underpin many algorithmic concepts.
    \item \textbf{“A Note on Two Problems in Connexion with Graphs” by Leonard Euler} - This seminal paper from 1736 is often considered the starting point of graph theory, introducing the famous Seven Bridges of Königsberg problem.
    \item \textbf{“The Shortest Path Problem” by Edsger W. Dijkstra} - Dijkstra's 1959 paper introduces his algorithm for finding the shortest path in a graph, a fundamental concept in graph theory.
\end{itemize}

\subsection{Online Tutorials and Courses}
Online tutorials and courses offer interactive and accessible ways to learn graph algorithms. These resources often include video lectures, coding exercises, and forums for discussion.

\begin{itemize}
    \item \textbf{Coursera’s “Algorithms” Specialization by Stanford University} - This series of courses covers a wide range of algorithmic topics, including several modules on graph algorithms.
    \item \textbf{edX’s “Algorithms and Data Structures” by IIT Bombay} - This course provides a comprehensive introduction to algorithms and data structures, with a strong emphasis on graph algorithms.
    \item \textbf{MIT OpenCourseWare’s “Introduction to Algorithms”} - MIT’s OCW offers free lecture notes, assignments, and exams for their algorithms course, which includes detailed coverage of graph algorithms.
    \item \textbf{Khan Academy’s “Algorithms” Course} - This course offers a user-friendly introduction to algorithms, including graph traversal techniques and shortest path algorithms.
    \item \textbf{LeetCode and HackerRank} - These platforms provide a plethora of problems and challenges on graph algorithms, allowing students to practice and refine their skills through hands-on coding.
\end{itemize}

\subsection{Open-Source Libraries and Tools for Graph Algorithm Implementation}
Implementing graph algorithms requires practical tools and libraries. Several open-source libraries can help you efficiently implement and test graph algorithms.

\begin{itemize}
    \item \textbf{NetworkX} - A Python library for the creation, manipulation, and study of the structure, dynamics, and functions of complex networks. NetworkX is particularly user-friendly for implementing graph algorithms.
    \begin{lstlisting}[language=Python, caption=Example of Dijkstra's Algorithm using NetworkX]
import networkx as nx

G = nx.Graph()
G.add_weighted_edges_from([(1, 2, 0.6), (1, 3, 0.2), (3, 4, 0.1), (4, 2, 0.7)])

# Compute shortest paths using Dijkstra's algorithm
shortest_path = nx.dijkstra_path(G, source=1, target=4)
print(shortest_path)
    \end{lstlisting}
    \FloatBarrier

    \item \textbf{Graph-tool} - An efficient Python module for manipulation and statistical analysis of graphs (networks). Graph-tool is designed for performance and scalability.
    \item \textbf{JGraphT} - A Java library that provides mathematical graph-theory objects and algorithms. It is highly versatile and can be used for a wide range of graph-related tasks.
    \item \textbf{Gephi} - An open-source network analysis and visualization software package written in Java. It allows users to interact with the representation, manipulation, and visualization of graphs.
    \item \textbf{Neo4j} - A highly scalable native graph database that leverages data relationships as first-class entities. Neo4j is particularly useful for applications that require efficient querying and manipulation of graph data.
    \item \textbf{igraph} - A collection of network analysis tools with interfaces to R, Python, and C/C++. It is particularly well-suited for large graphs and complex network analyses.
\end{itemize}

These resources will provide you with a robust toolkit for learning, implementing, and exploring graph algorithms. Whether you are just starting or looking to deepen your knowledge, these books, courses, and libraries offer valuable insights and practical skills.


\section{End of Chapter Exercises}
In this section, we provide a comprehensive set of exercises designed to reinforce and apply the concepts learned in this chapter on Graph Algorithm Techniques. These exercises are divided into two main categories: conceptual questions to reinforce learning and practical coding challenges to apply graph algorithm techniques. By completing these exercises, students will deepen their understanding of the material and gain practical experience in implementing and utilizing graph algorithms.

\subsection{Conceptual Questions to Reinforce Learning}
This subsection contains a series of conceptual questions aimed at reinforcing the theoretical understanding of graph algorithms. These questions are intended to test the students' grasp of key concepts, definitions, and properties related to graph theory and graph algorithms.

\begin{itemize}
    \item \textbf{Question 1:} Explain the difference between a directed and an undirected graph. Provide examples of real-world scenarios where each type would be applicable.
    \item \textbf{Question 2:} Define what a spanning tree is in the context of graph theory. How does it relate to a minimum spanning tree?
    \item \textbf{Question 3:} Describe Dijkstra's algorithm for finding the shortest path in a weighted graph. What are its time complexity and limitations?
    \item \textbf{Question 4:} Compare and contrast Depth-First Search (DFS) and Breadth-First Search (BFS) in terms of their methodologies and applications.
    \item \textbf{Question 5:} What is a bipartite graph? Provide a method to determine if a given graph is bipartite.
    \item \textbf{Question 6:} Discuss the concept of graph coloring. What is the chromatic number, and how can it be determined for a graph?
    \item \textbf{Question 7:} Explain the significance of the Bellman-Ford algorithm. In what scenarios is it preferable over Dijkstra's algorithm?
    \item \textbf{Question 8:} Describe the concept of network flow and the Ford-Fulkerson method for computing the maximum flow in a flow network.
\end{itemize}

\subsection{Practical Coding Challenges to Apply Graph Algorithm Techniques}
This subsection provides a set of practical coding challenges designed to apply the graph algorithm techniques discussed in this chapter. Each problem is accompanied by a detailed solution, including both the algorithmic description and the corresponding Python code implementation.

\begin{itemize}
    \item \textbf{Problem 1:} Implementing Dijkstra's Algorithm
    \item \textbf{Problem 2:} Detecting Cycles in a Graph using DFS
    \item \textbf{Problem 3:} Finding the Minimum Spanning Tree using Kruskal's Algorithm
    \item \textbf{Problem 4:} Determining if a Graph is Bipartite
\end{itemize}

\subsubsection*{Problem 1: Implementing Dijkstra's Algorithm}
Given a weighted graph, implement Dijkstra's algorithm to find the shortest path from a source node to all other nodes.

\begin{algorithm}
\caption{Dijkstra's Algorithm}
\begin{algorithmic}[1]
\State \textbf{Input:} Graph $G=(V, E)$ with edge weights $w$, source node $s$
\State \textbf{Output:} Shortest path distances from $s$ to all nodes in $V$
\State Initialize distance $d[v] \gets \infty$ for all $v \in V$
\State Set $d[s] \gets 0$
\State Initialize priority queue $Q$ with $(0, s)$
\While{$Q$ is not empty}
    \State Extract $(d[u], u)$ from $Q$
    \For{each neighbor $v$ of $u$}
        \If{$d[u] + w(u, v) < d[v]$}
            \State $d[v] \gets d[u] + w(u, v)$
            \State Insert $(d[v], v)$ into $Q$
        \EndIf
    \EndFor
\EndWhile
\State \Return $d$
\end{algorithmic}
\end{algorithm}

\FloatBarrier
\begin{lstlisting}[language=Python, caption=Python implementation of Dijkstra's Algorithm]
import heapq

def dijkstra(graph, start):
    # Initialize distances with infinity
    distances = {node: float('infinity') for node in graph}
    distances[start] = 0
    # Priority queue to store (distance, vertex) tuples
    priority_queue = [(0, start)]
    
    while priority_queue:
        current_distance, current_vertex = heapq.heappop(priority_queue)
        
        # Nodes can get added to the priority queue multiple times, we only
        # process a vertex the first time we remove it from the priority queue.
        if current_distance > distances[current_vertex]:
            continue
        
        for neighbor, weight in graph[current_vertex].items():
            distance = current_distance + weight
            
            # Only consider this new path if it's better
            if distance < distances[neighbor]:
                distances[neighbor] = distance
                heapq.heappush(priority_queue, (distance, neighbor))
    
    return distances

# Example usage
graph = {
    'A': {'B': 1, 'C': 4},
    'B': {'A': 1, 'C': 2, 'D': 5},
    'C': {'A': 4, 'B': 2, 'D': 1},
    'D': {'B': 5, 'C': 1}
}

start_node = 'A'
print(dijkstra(graph, start_node))
\end{lstlisting}
\FloatBarrier

\subsubsection*{Problem 2: Detecting Cycles in a Graph using DFS}
Implement a function to detect if a cycle exists in a directed graph using Depth-First Search (DFS).

\begin{algorithm}
\caption{Cycle Detection using DFS}
\begin{algorithmic}[1]
\State \textbf{Input:} Graph $G=(V, E)$
\State \textbf{Output:} Boolean indicating if a cycle exists
\State Initialize all vertices as unvisited
\State Initialize recursion stack $recStack$ as empty
\For{each vertex $v \in V$}
    \If{$v$ is unvisited}
        \If{\textsc{DFS-Visit}($v$)}
            \State \Return \textbf{true}
        \EndIf
    \EndIf
\EndFor
\State \Return \textbf{false}
\Function{DFS-Visit}{$v$}
    \State Mark $v$ as visited
    \State Add $v$ to $recStack$
    \For{each neighbor $u$ of $v$}
        \If{$u$ is unvisited and \textsc{DFS-Visit}($u$)}
            \State \Return \textbf{true}
        \ElsIf{$u$ is in $recStack$}
            \State \Return \textbf{true}
        \EndIf
    \EndFor
    \State Remove $v$ from $recStack$
    \State \Return \textbf{false}
\EndFunction
\end{algorithmic}
\end{algorithm}

\FloatBarrier
\begin{lstlisting}[language=Python, caption=Python implementation of Cycle Detection using DFS]
def has_cycle(graph):
    def dfs(vertex):
        visited.add(vertex)
        rec_stack.add(vertex)
        
        for neighbor in graph[vertex]:
            if neighbor not in visited:
                if dfs(neighbor):
                    return True
            elif neighbor in rec_stack:
                return True
        
        rec_stack.remove(vertex)
        return False
    
    visited = set()
    rec_stack = set()
    
    for node in graph:
        if node not in visited:
            if dfs(node):
                return True
    return False

# Example usage
graph = {
    'A': ['B'],
    'B': ['C'],
    'C': ['A'],
    'D': ['E'],
    'E': []
}

print(has_cycle(graph))  # Output: True
\end{lstlisting}
\FloatBarrier

\subsubsection*{Problem 3: Finding the Minimum Spanning Tree using Kruskal's Algorithm}
Implement Kruskal's algorithm to find the Minimum Spanning Tree (MST) of a given graph.

\begin{algorithm}
\caption{Kruskal's Algorithm}
\begin{algorithmic}[1]
\State \textbf{Input:} Graph $G=(V, E)$ with edge weights $w$
\State \textbf{Output:} Minimum spanning tree $T$
\State Sort edges $E$ by weight $w$
\State Initialize $T$ as empty
\State Initialize disjoint-set $D$ for $V$
\For{each edge $(u, v) \in E$ in ascending order of weight}
    \If{$u$ and $v$ are in different sets in $D$}
        \State Add $(u, v)$ to $T$
        \State Union the sets of $u$ and $v$ in $D$
    \EndIf
\EndFor
\State \Return $T$
\end{algorithmic}
\end{algorithm}

\FloatBarrier
\begin{lstlisting}[language=Python, caption=Python implementation of Kruskal's Algorithm]
class DisjointSet:
    def __init__(self, vertices):
        self.parent = {v: v for v in vertices}
        self.rank = {v: 0 for v in vertices}

    def find(self, vertex):
        if self.parent[vertex] != vertex:
            self.parent[vertex] = self.find(self.parent[vertex])
        return self.parent[vertex]

    def union(self, root1, root2):
        if self.rank[root1] > self.rank[root2]:
            self.parent[root2] = root1
        elif self.rank[root1] < self.rank[root2]:
            self.parent[root1] = root2
        else:
            self.parent[root2] = root1
            self.rank[root1] += 1

def kruskal(graph):
    edges = []
    for vertex in graph:
        for neighbor, weight in graph[vertex].items():
            edges.append((weight, vertex, neighbor))
    edges.sort()
    
    disjoint_set = DisjointSet(graph.keys())
    mst = []

    for weight, u, v in edges:
        root1 = disjoint_set.find(u)
        root2 = disjoint_set.find(v)
        if root1 != root2:
            mst.append((u, v, weight))
            disjoint_set.union(root1, root2)
    
    return mst

# Example usage
graph = {
    'A': {'B': 1, 'C': 3},
    'B': {'A': 1, 'C': 1, 'D': 6},
    'C': {'A': 3, 'B': 1, 'D': 5},
    'D': {'B': 6, 'C': 5}
}

print(kruskal(graph))
\end{lstlisting}
\FloatBarrier

\subsubsection*{Problem 4: Determining if a Graph is Bipartite}
Implement a function to check if a given graph is bipartite using BFS.

\begin{algorithm}
\caption{Bipartite Graph Check using BFS}
\begin{algorithmic}[1]
\State \textbf{Input:} Graph $G=(V, E)$
\State \textbf{Output:} Boolean indicating if the graph is bipartite
\State Initialize all vertices with no color
\For{each vertex $v \in V$}
    \If{$v$ is uncolored}
        \State Color $v$ with color 1
        \If{\textsc{BFS-Check}($v$) returns false}
            \State \Return \textbf{false}
        \EndIf
    \EndIf
\EndFor
\State \Return \textbf{true}
\Function{BFS-Check}{$v$}
    \State Initialize queue $Q$ with $v$
    \While{$Q$ is not empty}
        \State Dequeue $u$ from $Q$
        \For{each neighbor $w$ of $u$}
            \If{$w$ is uncolored}
                \State Color $w$ with the opposite color of $u$
                \State Enqueue $w$ into $Q$
            \ElsIf{$w$ has the same color as $u$}
                \State \Return \textbf{false}
            \EndIf
        \EndFor
    \EndWhile
    \State \Return \textbf{true}
\EndFunction
\end{algorithmic}
\end{algorithm}

\FloatBarrier
\begin{lstlisting}[language=Python, caption=Python implementation of Bipartite Check using BFS]
from collections import deque

def is_bipartite(graph):
    color = {}
    
    def bfs(start):
        queue = deque([start])
        color[start] = 0
        
        while queue:
            node = queue.popleft()
            for neighbor in graph[node]:
                if neighbor not in color:
                    color[neighbor] = 1 - color[node]
                    queue.append(neighbor)
                elif color[neighbor] == color[node]:
                    return False
        return True
    
    for node in graph:
        if node not in color:
            if not bfs(node):
                return False
    return True

# Example usage
graph = {
    'A': ['B', 'C'],
    'B': ['A', 'D'],
    'C': ['A', 'D'],
    'D': ['B', 'C']
}

print(is_bipartite(graph))  # Output: True
\end{lstlisting}
\FloatBarrier
\end{document}